% ============================================================
%  PIXEL-2026 — Technical & Research Background Document
%  Target Venue : AAAI-2027  (+ IEEE TMTT / IEEE TNNLS)
%  Status       : Active Research Blueprint  |  v2.0  |  May 2026
% ============================================================

\documentclass[12pt,a4paper]{article}

% ---------- encoding & language ----------
\usepackage[T1]{fontenc}
\usepackage[utf8]{inputenc}
\usepackage[english]{babel}

% ---------- page geometry ----------
\usepackage[margin=1in, top=1.1in, bottom=1.1in]{geometry}

% ---------- math ----------
\usepackage{amsmath, amssymb, amsthm, mathtools}
\usepackage{bm}          % bold math symbols
\usepackage{bbm}         % blackboard-bold 1

% ---------- fonts ----------
\usepackage{lmodern}
\usepackage{microtype}

% ---------- graphics & colour ----------
\usepackage{graphicx}
\usepackage[dvipsnames,table]{xcolor}
\definecolor{pixelblue}{RGB}{0,74,173}
\definecolor{pixelgray}{RGB}{245,245,248}
\definecolor{warnred}{RGB}{180,20,20}
\definecolor{okgreen}{RGB}{0,120,60}

% ---------- tables ----------
\usepackage{booktabs}
\usepackage{tabularx}
\usepackage{multirow}
\usepackage{array}

% ---------- lists ----------
\usepackage{enumitem}
\setlist[itemize]{leftmargin=1.5em, itemsep=2pt}
\setlist[enumerate]{leftmargin=1.5em, itemsep=2pt}

% ---------- code / algorithms ----------
\usepackage{listings}
\lstset{
  backgroundcolor=\color{pixelgray},
  basicstyle=\ttfamily\small,
  frame=single,
  rulecolor=\color{gray!40},
  breaklines=true,
  numbers=none
}
\usepackage[ruled,vlined,linesnumbered]{algorithm2e}

% ---------- hyperlinks ----------
\usepackage[
  colorlinks=true,
  linkcolor=pixelblue,
  citecolor=pixelblue,
  urlcolor=pixelblue,
  pdftitle={PIXEL-2026 Technical Background},
  pdfauthor={PIXEL Research Team}
]{hyperref}

% ---------- section styling ----------
\usepackage{titlesec}
\titleformat{\section}{\large\bfseries\color{pixelblue}}{\thesection.}{0.6em}{}[\titlerule]
\titleformat{\subsection}{\normalsize\bfseries\color{pixelblue!80!black}}{\thesubsection}{0.5em}{}
\titleformat{\subsubsection}{\normalsize\bfseries}{\thesubsubsection}{0.5em}{}

% ---------- boxes for notes ----------
\usepackage[most]{tcolorbox}
\tcbset{
  notebox/.style={
    colback=pixelgray, colframe=pixelblue!60,
    fonttitle=\bfseries\small, boxrule=0.5pt,
    arc=3pt, left=6pt, right=6pt, top=4pt, bottom=4pt
  },
  warnbox/.style={
    colback=red!5, colframe=warnred!70,
    fonttitle=\bfseries\small, boxrule=0.5pt,
    arc=3pt, left=6pt, right=6pt, top=4pt, bottom=4pt
  },
  fixbox/.style={
    colback=green!5, colframe=okgreen!70,
    fonttitle=\bfseries\small, boxrule=0.5pt,
    arc=3pt, left=6pt, right=6pt, top=4pt, bottom=4pt
  }
}

% ---------- misc ----------
\usepackage{caption}
\captionsetup{font=small, labelfont=bf}
\usepackage{float}
\usepackage{parskip}
\setlength{\parskip}{6pt}

% ============================================================
%  CUSTOM COMMANDS
% ============================================================
\newcommand{\FEM}{\mathcal{F}_{\mathrm{EM}}}
\newcommand{\xtilde}{\tilde{x}}
\newcommand{\xhat}{\hat{x}}
\newcommand{\ystar}{\mathbf{y}^{*}}
\newcommand{\yhat}{\hat{\mathbf{y}}}
\newcommand{\Smat}{\mathbf{S}}
\newcommand{\cy}{\mathbf{c}_y}
\newcommand{\HW}{H \times W}
\newcommand{\tthresh}{t_{\mathrm{thresh}}}
\newcommand{\alphamax}{\alpha_{\max}}
\newcommand{\sigmahat}{\hat{\sigma}}
\newcommand{\Real}{\operatorname{Re}}
\newcommand{\Imag}{\operatorname{Im}}
\newcommand{\Hilbert}{\mathcal{H}}
\newcommand{\NOTE}[1]{\begin{tcolorbox}[notebox]#1\end{tcolorbox}}
\newcommand{\WARN}[1]{\begin{tcolorbox}[warnbox, title=V1 Error — Corrected in V2]#1\end{tcolorbox}}
\newcommand{\FIX}[1]{\begin{tcolorbox}[fixbox, title=V2 Fix]#1\end{tcolorbox}}

% ============================================================
%  THEOREMS / DEFINITIONS
% ============================================================
\theoremstyle{definition}
\newtheorem{definition}{Definition}[section]
\newtheorem{hypothesis}{Hypothesis}[section]
\theoremstyle{plain}
\newtheorem{proposition}{Proposition}[section]
\theoremstyle{remark}
\newtheorem{remark}{Remark}[section]

% ============================================================
\begin{document}
% ============================================================

% ----------------------------------------------------------
%  TITLE PAGE
% ----------------------------------------------------------
\begin{titlepage}
  \centering
  \vspace*{1.5cm}

  {\Huge\bfseries\color{pixelblue}
    Physics-Constrained Probabilistic\\[6pt]
    Topology Synthesis for Inverse\\[6pt]
    Electromagnetic RF/IC Design\\[10pt]
  }

  \vspace{0.5cm}
  {\large\color{gray} \textit{PIXEL-2026 — Technical \& Research Background Document}}

  \vspace{1cm}
  \rule{0.8\linewidth}{1pt}

  \vspace{0.5cm}
  {\normalsize
    \begin{tabular}{ll}
      \textbf{Document Version:} & 2.0 (Validated \& Error-Corrected) \\[4pt]
      \textbf{Date:}             & May 2026 \\[4pt]
      \textbf{Target Venue:}     & AAAI-2027 \\[4pt]
      \textbf{Concurrent:}       & IEEE TMTT / IEEE TNNLS \\[4pt]
      \textbf{Status:}           & Active Research Blueprint \\
    \end{tabular}
  }

  \vspace{1cm}
  \rule{0.8\linewidth}{0.5pt}

  \vspace{0.8cm}
  \begin{tcolorbox}[notebox, width=0.85\linewidth]
    \textbf{Document Purpose.}
    This is the authoritative master research context for the PIXEL-2026 project.
    It covers the full problem from first principles—physics, mathematics, AI
    methodology, engineering constraints, novelty analysis, vulnerability assessment,
    and AAAI positioning.  Every section has been validated, corrected, and enriched
    from the v1 brainstorm.  All mathematical errors present in v1 have been
    identified, documented, and fixed.
  \end{tcolorbox}

  \vfill
  {\small\color{gray}
    \textit{Hare Krishna.  All knowledge is offered at the lotus feet of the Supreme.}
  }
\end{titlepage}

% ----------------------------------------------------------
%  ABSTRACT
% ----------------------------------------------------------
\begin{abstract}
\noindent
We present \textbf{PIXEL} (\textbf{P}hysics-guided d\textbf{I}screte di\textbf{X}fusion for
\textbf{E}M \textbf{L}ayout synthesis), a novel framework for inverse electromagnetic
(EM) design of RF/IC structures.
Given a target spectral specification expressed as S-parameters, PIXEL generates
fabrication-ready, Maxwell-consistent binary conductor layouts on a discretised substrate
grid.

The framework rests on three interlocking scientific contributions:
\begin{enumerate}
  \item A \textbf{discrete diffusion generative model} (D3PM) over the binary EM layout space,
        correctly formulated for ternary-state $\{0, 1, \textsc{mask}\}$ absorption-based corruption.
  \item \textbf{Uncertainty-weighted differentiable physics guidance} that steers the reverse
        denoising trajectory toward spectral objectives using a calibrated ensemble surrogate of
        the full-wave EM forward map.
  \item \textbf{Differentiable connectivity and manufacturability constraints} embedded
        inside the generation process—not applied as post-hoc corrections.
\end{enumerate}

All v1 mathematical errors (continuous CFG applied to discrete space; gradient guidance on
discrete tokens; missing Kramers--Kronig causality; incomplete passivity constraint) are
formally corrected in this document.  The framework targets AAAI-2027 and concurrent
IEEE TMTT submission.
\end{abstract}

\newpage
\tableofcontents
\newpage

% ============================================================
\section{Executive Summary}
\label{sec:exec}
% ============================================================

The inverse electromagnetic design problem asks:
\emph{given a desired radio-frequency response, what physical layout of conductors on a
dielectric substrate produces that response?}
This problem is \textbf{ill-posed}, \textbf{combinatorially hard}, and of enormous
engineering relevance—every modern RF chip, filter, coupler, and antenna requires solving
some variant of it.

Classical approaches either enumerate topology candidates using expensive full-wave
simulations or restrict the search to pre-defined circuit templates.
Neural inverse regressors learn a fast mapping from specification to layout, but because
the mapping is one-to-many, deterministic networks collapse to physically meaningless
``average'' structures.

\medskip
PIXEL resolves this by treating inverse EM design as \textbf{constrained probabilistic
topology generation}.
The key ideas are:

\begin{itemize}
  \item \textbf{Binary discrete diffusion.}
        EM layouts live in $\{0,1\}^{H \times W}$—a discrete, binary space.
        Gaussian (continuous) diffusion is the wrong model.
        D3PM with absorbing-state masking is the right one.
  \item \textbf{Physics guidance via differentiable surrogate.}
        A fast CNN ensemble approximates the full-wave solver.
        Its gradients steer the generative trajectory toward the target S-parameters.
  \item \textbf{Topology and manufacturability as differentiable losses.}
        A learned connectivity discriminator and differentiable DRC losses prevent the
        generation of fabrication-invalid structures.
\end{itemize}

\begin{tcolorbox}[notebox, title=Critical Corrections from V1]
\begin{itemize}
  \item CFG formula copied verbatim from continuous diffusion — \textbf{wrong for discrete
        diffusion} (fixed in Section~\ref{sec:cfg}).
  \item Physics guidance gradient applied to discrete tokens — \textbf{undefined} (fixed in
        Section~\ref{sec:guidance}).
  \item Kramers--Kronig causality constraints entirely missing (added in
        Section~\ref{sec:physics}).
  \item Multi-port passivity constraint was scalar-incomplete (full matrix form in
        Section~\ref{sec:physics}).
\end{itemize}
\end{tcolorbox}

% ============================================================
\section{Problem Statement}
\label{sec:problem}
% ============================================================

\subsection{What We Are Trying to Solve}

Imagine you are designing a bandpass filter for a 5G chip.
You know exactly what you want: signals between 3.5~GHz and 3.7~GHz should pass through
with less than 0.5~dB loss; everything outside that band should be suppressed by at least
20~dB.
What you \emph{do not} know is the physical shape of the copper trace on the PCB that will
achieve this.
That is the inverse EM design problem—going from \emph{what you want} to \emph{what you
build}.

\subsection{Formal Problem Definition}
\label{sec:problem:formal}

\begin{definition}[Inverse EM Design Problem]
\textbf{Given:}
\begin{itemize}
  \item Target S-parameter specification $\ystar \in \mathbb{R}^{d}$, encoding
        $|S_{11}(f)|$, $|S_{21}(f)|$, phase, and group delay across $N_f$ frequency samples.
  \item Substrate specification $(\varepsilon_r,\, \tan\delta,\, h)$: relative permittivity,
        loss tangent, substrate height.
  \item Fabrication specification $(w_{\min},\, s_{\min},\, \Delta_{\mathrm{DRC}})$: minimum
        trace width, minimum spacing, and the full DRC rule set.
  \item Operating frequency range $[f_{\min},\, f_{\max}]$.
\end{itemize}

\textbf{Generate:}
\begin{itemize}
  \item A physical binary layout $x \in \{0,1\}^{\HW}$, where $1 =$ conductor and
        $0 =$ dielectric/void.
  \item Such that $\FEM(x) \approx \ystar$ (Maxwell-consistent spectral behavior).
  \item And $x$ is topologically valid: connected, with port-accessible conducting paths.
  \item And $x$ satisfies all manufacturing DRC rules.
\end{itemize}
\end{definition}

\subsection{Why the Problem Is Hard}
\label{sec:problem:hard}

\begin{table}[H]
\centering
\caption{Sources of hardness in the inverse EM design problem.}
\label{tab:hardness}
\begin{tabular}{lp{7.5cm}}
\toprule
\textbf{Property} & \textbf{Consequence} \\
\midrule
Binary domain $\{0,1\}^{\HW}$ &
  Search space $2^{HW}$; for $H=W=15$: $2^{225} \approx 10^{67}$ possibilities. \\
\addlinespace
Topology sensitivity &
  A single pixel flip can cause a discontinuous jump in the S-parameter response. \\
\addlinespace
Many-to-one forward map &
  Multiple distinct layouts produce nearly identical spectra (ill-posed inverse). \\
\addlinespace
Non-convex forward physics &
  Global optimisation of $\FEM$ is computationally intractable. \\
\addlinespace
Fabrication constraints &
  Reduces the valid search space but adds further combinatorial structure. \\
\bottomrule
\end{tabular}
\end{table}

The problem is \textbf{ill-posed} in the sense of Hadamard:
\begin{itemize}
  \item \textbf{Non-uniqueness.}  For any $\ystar$, multiple layouts $x$ satisfy
        $\FEM(x) \approx \ystar$.
  \item \textbf{Non-existence.}  Some $\ystar$ may not be physically realizable at a
        given resolution.
  \item \textbf{Instability.}  Small perturbations in $\ystar$ may yield topologically
        different layouts.
\end{itemize}

These three properties together mandate a \textbf{probabilistic} approach: instead of
finding a single layout, we learn the full conditional distribution
$P(x \mid \ystar)$ over valid layouts.

\subsection{Why Existing Approaches Fail}
\label{sec:problem:existing}

\paragraph{Deterministic CNN/MLP regressors.}
Output a single $x$ for each $\ystar$.
For ill-posed mappings, training collapses to the conditional mean—a blurry,
topologically invalid, physically nonsensical average over the solution manifold.

\paragraph{Genetic algorithms / evolutionary methods.}
Require $\mathcal{O}(10^3$--$10^5)$ EM forward evaluations per design.
Non-differentiable and poor at scaling with layout complexity.
No concept of a learned design manifold.

\paragraph{Adjoint topology optimisation.}
Requires a continuous material density (not discrete binary topology).
Produces grayscale designs that must be projected to binary—introducing spectral
artefacts and manufacturability violations.

\paragraph{Standard continuous diffusion (DDPM/DDIM).}
Assumes Gaussian noise corruption.
Binary EM layouts violate this assumption: intermediate diffusion states are
continuous-valued with no physical interpretation, and the learned prior is wrong.

\NOTE{\textbf{Key takeaway.}
Each existing method fails either because it ignores the probabilistic nature of the
problem, assumes continuous variables when the domain is discrete, or cannot enforce
topological validity during generation.
PIXEL addresses all three failures simultaneously.}

% ============================================================
\section{Physics Foundations and First Principles}
\label{sec:physics}
% ============================================================

\subsection{Why Physics Matters in an AI Paper}

Before building any neural network, we must understand \emph{what physics governs the
thing we are designing}.
EM structures obey Maxwell's equations—these are not soft guidelines but hard laws.
Any generated layout that violates them is physically meaningless, regardless of how
good the neural network loss looks.
This section establishes the physical ground truth that every component of PIXEL must
respect.

\subsection{Governing Equations (Maxwell in Frequency Domain)}
\label{sec:maxwell}

At a single frequency $f$ (with $\omega = 2\pi f$), electromagnetic behaviour is
described by the time-harmonic (phasor) form of Maxwell's equations:

\begin{align}
  \nabla \times \mathbf{E} &= -j\omega\mu\mathbf{H} \label{eq:faraday}\\
  \nabla \times \mathbf{H} &= \mathbf{J} + j\omega\varepsilon\mathbf{E} \label{eq:ampere}\\
  \nabla \cdot \mathbf{D}  &= \rho_f \label{eq:gauss_e}\\
  \nabla \cdot \mathbf{B}  &= 0 \label{eq:gauss_m}
\end{align}

The binary layout $x \in \{0,1\}^{\HW}$ determines the \emph{spatial permittivity
distribution}:

\begin{equation}
  \varepsilon(\mathbf{r}) =
  \begin{cases}
    \varepsilon_0\!\left(1 - \dfrac{j\sigma}{\omega\varepsilon_0}\right) &
      \text{if pixel $(i,j)$ containing $\mathbf{r}$ has } x_{ij} = 1
      \quad (\text{conductor}) \\[6pt]
    \varepsilon_0 \varepsilon_r &
      \text{if } x_{ij} = 0 \quad (\text{dielectric})
  \end{cases}
  \label{eq:permittivity}
\end{equation}

\textbf{The inverse design task is therefore:} find the binary map $x$ (equivalently,
find $\varepsilon(\mathbf{r})$) such that the resulting S-parameters match $\ystar$.

\subsection{Scattering Matrix Theory — Complete Formulation}
\label{sec:smatrix}

For an $N$-port passive microwave network, the scattering matrix
$\Smat \in \mathbb{C}^{N \times N}$ must satisfy three physical laws:

\paragraph{Passivity (power conservation).}
\begin{equation}
  \Smat^\dagger \Smat \preceq \mathbf{I}
  \label{eq:passivity}
\end{equation}
The notation $\preceq \mathbf{I}$ means the matrix $(\mathbf{I} - \Smat^\dagger\Smat)$
is positive semi-definite—no more power can come out of a passive network than goes in.
For a 2-port this reduces to the familiar $|S_{11}|^2 + |S_{21}|^2 \leq 1$,
but the \emph{full matrix form must be used} in the training loss (the scalar form is
incorrect for $N > 2$).

\paragraph{Reciprocity (for non-magnetic media).}
\begin{equation}
  \Smat^T = \Smat \quad \Longrightarrow \quad S_{ij} = S_{ji}
  \label{eq:reciprocity}
\end{equation}
In particular $S_{12} = S_{21}$ for 2-port structures; exploiting this halves the
output dimension of the surrogate.

\paragraph{Losslessness (ideally lossless structures only).}
\begin{equation}
  \Smat^\dagger \Smat = \mathbf{I}
  \label{eq:lossless}
\end{equation}
Real microstrip substrates have dielectric and conductor losses, so losslessness is
unrealistic.  The passivity inequality~\eqref{eq:passivity} is the correct soft constraint.

\subsection{Causality and Kramers--Kronig Relations}
\label{sec:kk}

\begin{tcolorbox}[warnbox, title=V1 Gap — Critical Missing Physics]
Version~1 contained no causality constraint.
A surrogate trained without it can learn physically impossible S-parameter responses
where the output \emph{precedes} the input in time.
\end{tcolorbox}

Physical S-parameters must be \textbf{causal}: the response cannot precede the
excitation.
This is enforced by the \textbf{Kramers--Kronig (KK) relations}, which link the real
and imaginary parts of any linear causal transfer function:

\begin{align}
  \Real[S_{ij}(f)] &=
    \frac{2}{\pi}\,\mathcal{P}\!\int_0^\infty
      \frac{f'\,\Imag[S_{ij}(f')]}{f'^2 - f^2}\,df'
  \label{eq:kk_real}\\[4pt]
  \Imag[S_{ij}(f)] &=
    -\frac{2}{\pi}\,\mathcal{P}\!\int_0^\infty
      \frac{f\,\Real[S_{ij}(f')]}{f'^2 - f^2}\,df'
  \label{eq:kk_imag}
\end{align}

where $\mathcal{P}$ denotes the Cauchy principal value integral.

\textbf{Critical implementation consequence.}
Training a surrogate to independently regress magnitude and phase as uncorrelated targets
violates causality.
The fix: enforce the KK constraint via a loss term during surrogate training.

\begin{equation}
  \mathcal{L}_{KK} =
    \left\|\Real[\hat{\Smat}] - \Hilbert\!\left\{\Imag[\hat{\Smat}]\right\}\right\|^2
  \label{eq:lkk}
\end{equation}

where $\Hilbert$ is the discrete Hilbert transform operator (computed via FFT in practice).

\subsection{Transmission Line Theory — Grounding the Dataset}
\label{sec:tl}

Classical RF design uses the \textbf{telegrapher's equations}:

\begin{align}
  \frac{dV(z)}{dz} &= -(R + j\omega L)\,I(z) \\
  \frac{dI(z)}{dz} &= -(G + j\omega C)\,V(z)
\end{align}

with propagation constant
$\gamma = \sqrt{(R+j\omega L)(G+j\omega C)}$.
For a lossless line: $\gamma = j\beta = j\omega\sqrt{LC}$.

Every \emph{primitive structure} in our procedural dataset is the pixel-discretised
realisation of a classical transmission-line element.
This is not an arbitrary design choice—it grounds the dataset in known physics:

\begin{table}[H]
\centering
\caption{Mapping from classical RF elements to pixel-grid primitives.}
\label{tab:primitives}
\begin{tabular}{lll}
\toprule
\textbf{Classical Element} & \textbf{Pixel Primitive} & \textbf{Physical Behaviour} \\
\midrule
Microstrip line (length $\ell$)   & Row of connected pixels      & Transmission with phase shift $\beta\ell$ \\
$\lambda/4$ shunt stub            & Vertical stub at midpoint    & Bandstop at design frequency $f_0$ \\
$\lambda/2$ resonator             & Isolated line segment        & Bandpass resonance \\
Coupled lines                     & Two parallel adjacent traces & Directional coupling \\
Interdigital structure            & Alternating finger pixels    & Wideband bandpass filter \\
Split-ring resonator (SRR)        & Broken-ring pixel pattern    & Negative effective permeability \\
\bottomrule
\end{tabular}
\end{table}

\subsection{Effective Medium Theory and Pixel Discretisation}
\label{sec:emt}

At sub-wavelength scales, the binary pixel assumption is valid only when the pixel size
is much smaller than the effective wavelength.
The Maxwell--Garnett effective medium theory gives:

\begin{equation}
  \varepsilon_{\mathrm{eff}} =
    \varepsilon_h\,
    \frac{\varepsilon_i(1+2f) + 2\varepsilon_h(1-f)}
         {\varepsilon_i(1-f)  + \varepsilon_h(2+f)}
  \label{eq:maxwell_garnett}
\end{equation}

where $f$ is the conductor fill fraction per pixel, $\varepsilon_i$ is conductor
permittivity, and $\varepsilon_h$ is host (substrate) permittivity.

For the binary pixel assumption to hold, the pixel size must satisfy:

\begin{equation}
  \Delta_{\mathrm{pixel}} \ll \frac{\lambda_{\mathrm{eff}}}{10},
  \qquad
  \lambda_{\mathrm{eff}} = \frac{c}{f_{\max}\sqrt{\varepsilon_{r,\mathrm{eff}}}}
  \label{eq:pixel_validity}
\end{equation}

This constraint \emph{directly couples} the grid resolution, the frequency range, and
the substrate properties.

\subsection{Resolution–Frequency Co-Design}
\label{sec:resfcodesign}

\begin{tcolorbox}[warnbox, title=V1 Gap — Now Resolved]
Version~1 chose $H = W = 15$ without physical justification.
The co-design table below shows the valid frequency range for different physical domain
sizes at $15\times15$ resolution on Rogers~4003C ($\varepsilon_r = 3.5$).
\end{tcolorbox}

\begin{table}[H]
\centering
\caption{Valid operating frequency for $15\!\times\!15$ pixel grid
         (Rogers~4003C, $\varepsilon_r=3.5$).
         Pixel validity criterion: $\Delta < \lambda_{\mathrm{eff}}/10$.}
\label{tab:res_freq}
\begin{tabular}{ccc}
\toprule
\textbf{Physical domain $L$} & \textbf{Pixel size $\Delta = L/15$} & \textbf{Valid $f_{\max}$} \\
\midrule
20~mm & 1.33~mm & $\approx$4.5~GHz \\
10~mm & 0.67~mm & $\approx$9~GHz   \\
 5~mm & 0.33~mm & $\approx$18~GHz  \\
\bottomrule
\end{tabular}
\end{table}

\textbf{Recommendation.}
Constrain the physical domain to $5$--$10$~mm at $15\times15$ resolution.
This covers the range 2--18~GHz, which includes most RF/microwave applications
(cellular, WiFi, Bluetooth, satellite bands).

\NOTE{\textbf{Scaling plan.}
The $15\times15$ grid is a valid proof-of-concept.
The architecture naturally extends to $32\times32$ (covering mmWave up to $\sim$40~GHz)
with a larger dataset and higher-resolution simulations.
A scaling experiment at $32\times32$ is included in the evaluation plan.}

% ============================================================
\section{Mathematical Formulation}
\label{sec:math}
% ============================================================

\subsection{Layout Space}

\begin{equation}
  x \in \{0,1\}^{\HW}, \quad H = W = 15
  \label{eq:layout_space}
\end{equation}

A pixel value of $1$ denotes conductor; $0$ denotes dielectric.
For training purposes, a \textbf{continuous relaxation} is required:

\begin{equation}
  \xtilde \in [0,1]^{\HW}
  \label{eq:relaxed}
\end{equation}

with annealed binarisation during training (detailed in Section~\ref{sec:binarisation}).

\subsection{Spectral Specification Space}

\begin{equation}
  \mathbf{y} \in \mathbb{R}^{d}, \quad
  d = N_{\mathrm{port}}^2 \times N_f \times 2
  \label{eq:spec_space}
\end{equation}

The factor of~2 accounts for magnitude and phase (or equivalently, real and imaginary
parts).
For a 2-port with $N_f = 100$ frequency samples:
$d = 4 \times 100 \times 2 = 800$.
Exploiting reciprocity ($S_{12} = S_{21}$) reduces to $d = 2 \times N_f \times 2 = 400$
in practice.

\subsection{Forward Electromagnetic Mapping}

\begin{equation}
  \mathbf{y} = \FEM(x)
  \label{eq:forward}
\end{equation}

$\FEM$ denotes the full-wave EM solver (OpenEMS FDTD in this work).
Key properties:

\begin{itemize}
  \item \textbf{Nonlinear} — S-parameters are a nonlinear function of topology.
  \item \textbf{Non-differentiable w.r.t.\ discrete $x$} — a pixel flip is a discontinuous event.
  \item \textbf{Expensive} — each evaluation takes 1--10~minutes (frequency-swept FDTD).
  \item \textbf{Physics-complete} — satisfies Maxwell's equations by construction
        (within discretisation error).
\end{itemize}

\subsection{Inverse Design Objective (Extended)}
\label{sec:objective}

The formal optimisation objective is:

\begin{equation}
  x^* = \arg\min_{x \in \mathcal{X}_{\mathrm{valid}}}
    \mathcal{L}_{\mathrm{spec}}\!\left(\FEM(x),\, \ystar\right)
    + \lambda_{KK}\,\mathcal{L}_{KK}(x)
    + \lambda_{\mathrm{pass}}\,\mathcal{L}_{\mathrm{pass}}(x)
  \label{eq:objective}
\end{equation}

where $\mathcal{X}_{\mathrm{valid}}$ is the feasible set requiring:
topological validity (connected, port-accessible),
manufacturing DRC compliance, and physical realizability (no floating islands).

Direct optimisation over $\mathcal{X}_{\mathrm{valid}}$ is intractable.
We instead learn the conditional distribution:

\begin{equation}
  P_\theta(x \mid \mathbf{y})
    \propto P_\theta(x)\cdot P(\mathbf{y} \mid x)
    \cdot \mathbbm{1}[x \in \mathcal{X}_{\mathrm{valid}}]
  \label{eq:posterior}
\end{equation}

\subsection{Spectral Loss Function}

\begin{equation}
  \mathcal{L}_{\mathrm{spec}}(\yhat, \ystar)
  =
  \underbrace{
    \frac{1}{N_f}\sum_f \bigl(|S_{11}(f)| - |S_{11}^*(f)|\bigr)^2
  }_{\text{return loss}}
  +
  \underbrace{
    \frac{1}{N_f}\sum_f \bigl(|S_{21}(f)| - |S_{21}^*(f)|\bigr)^2
  }_{\text{insertion loss}}
  \label{eq:lspec}
\end{equation}

For filter designs, band-specific losses are added:

\begin{equation}
  \mathcal{L}_{\mathrm{band}}
  =
  \frac{1}{|\mathcal{F}_{\mathrm{pass}}|}\sum_{f \in \mathcal{F}_{\mathrm{pass}}}
    |S_{21}(f)|^2
  +
  \frac{1}{|\mathcal{F}_{\mathrm{stop}}|}\sum_{f \in \mathcal{F}_{\mathrm{stop}}}
    \bigl(|S_{21}(f)| - S_{\mathrm{stop}}^*\bigr)^2
  \label{eq:lband}
\end{equation}

\subsection{Surrogate Physics Loss (Full Form)}

The surrogate is trained with physics-consistency regularisers:

\begin{align}
  \mathcal{L}_{\mathrm{surrogate}}
  &= \underbrace{\|\yhat - \mathbf{y}\|_2^2}_{\mathrm{MSE}}
   + \lambda_{\mathrm{pass}} \underbrace{\mathcal{L}_{\mathrm{pass}}}_{\mathrm{passivity}}
   + \lambda_{\mathrm{recip}} \underbrace{\mathcal{L}_{\mathrm{recip}}}_{\mathrm{reciprocity}}
   + \lambda_{KK} \underbrace{\mathcal{L}_{KK}}_{\mathrm{causality}}
   + \lambda_{\mathrm{smooth}} \underbrace{\mathcal{L}_{\mathrm{smooth}}}_{\mathrm{smoothness}}
  \label{eq:lsurr}
\end{align}

\paragraph{Passivity loss (full matrix form).}
\begin{equation}
  \mathcal{L}_{\mathrm{pass}}
  = \frac{1}{N_f}\sum_f
    \max\!\bigl(0,\; \lambda_{\max}(\hat{\Smat}^\dagger(f)\hat{\Smat}(f)) - 1\bigr)^2
  \label{eq:lpass}
\end{equation}

where $\lambda_{\max}(\cdot)$ is the maximum eigenvalue.

\paragraph{Reciprocity loss.}
\begin{equation}
  \mathcal{L}_{\mathrm{recip}}
  = \frac{1}{N_f}\sum_f \|\hat{\Smat}(f) - \hat{\Smat}^T(f)\|_F^2
  \label{eq:lrecip}
\end{equation}

\paragraph{Spectral smoothness.}
\begin{equation}
  \mathcal{L}_{\mathrm{smooth}}
  = \frac{1}{N_f - 1}\sum_f
    \|\hat{\Smat}(f+\Delta f) - \hat{\Smat}(f)\|_F^2
  \label{eq:lsmooth}
\end{equation}

This suppresses spurious high-frequency artefacts in the predicted spectrum.

% ============================================================
\section{Core Scientific Hypothesis}
\label{sec:hypothesis}
% ============================================================

\begin{hypothesis}[PIXEL Central Claim]
\label{hyp:main}
Inverse EM design can be formulated as constrained probabilistic topology synthesis
over a \emph{learned manifold} of physically plausible electromagnetic structures,
where differentiable surrogate physics and topological constraints jointly guide
discrete generative sampling toward desired spectral objectives.
\end{hypothesis}

This hypothesis has three parts, each independently supported.

\subsection{Part 1 — The EM Layout Manifold Is Low-Dimensional}

Out of $2^{225} \approx 10^{67}$ possible $15\times15$ binary images, only an
astronomically small fraction are:
\begin{itemize}
  \item physically connected (most random binary images have isolated conductors),
  \item non-trivially resonant (structured resonances, controlled transmission), and
  \item manufacturable (satisfy DRC rules).
\end{itemize}

This tiny fraction forms a \textbf{low-dimensional manifold} within the binary space.
Learning this manifold is precisely what the generative prior must accomplish—and this
is strongly supported by all classical physics of RF structures.

\subsection{Part 2 — A Generative Prior Can Learn This Manifold}

Successful precedents in analogous scientific domains:
\begin{itemize}
  \item \textbf{Protein structure generation:} Discrete amino-acid sequences + structural constraints → folded proteins.
  \item \textbf{Molecule generation:} Graph-based discrete structures with physical validity constraints.
  \item \textbf{Crystal structure generation:} Periodic discrete lattices with symmetry constraints.
\end{itemize}

The EM case is arguably simpler: physical constraints are either entirely local
(topology, minimum feature size) or global (spectral behaviour), exactly matching the
inductive bias of graph-conditioned discrete diffusion.

\subsection{Part 3 — Differentiable Surrogate Physics Can Guide Sampling}

\begin{itemize}
  \item Analogous to \textbf{classifier guidance} in image diffusion
        (Dhariwal \& Nichol, 2021), but with a physics surrogate replacing a classifier.
  \item The ``class label'' is now a continuous spectral specification.
  \item The ``classifier'' is the EM surrogate model.
  \item This combination has direct precedent in molecule property-guided diffusion
        (DiffSBDD, Schneuing et al., 2022).
\end{itemize}

% ============================================================
\section{Novel Contributions — Detailed and Defended}
\label{sec:contributions}
% ============================================================

\subsection{Contribution 1: Discrete Physics-Guided Generative Framework}

\paragraph{What it is.}
A complete end-to-end system combining discrete diffusion (D3PM) over binary EM layouts
with differentiable physics guidance during reverse sampling.

\paragraph{Why it is novel.}
No prior work combines:
\begin{enumerate}
  \item Binary discrete diffusion (not continuous) for EM structures,
  \item Physics surrogate gradient guidance in the discrete reverse process, and
  \item Topology constraints embedded in the generative objective.
\end{enumerate}

\paragraph{AAAI argument.}
The paper makes a methodological contribution to \emph{constrained discrete generative
modelling} with domain application in EM design.
The methodology is novel independent of the application domain.

\subsection{Contribution 2: Uncertainty-Weighted Physics Guidance}

\paragraph{What it is.}
The physics guidance step size is modulated by the ensemble surrogate's uncertainty:

\begin{equation}
  \alpha_t \propto \frac{1}{\sigmahat(\xhat_0) + \epsilon}
  \label{eq:uncertainty_guidance}
\end{equation}

where $\sigmahat(\xhat_0)$ is the ensemble surrogate variance at the current
(partially denoised) predicted layout.

\paragraph{Why it is novel.}
Standard guidance uses a fixed or timestep-decayed step size.
Uncertainty-weighted guidance provides a principled mechanism to:
\begin{itemize}
  \item Apply \emph{strong} guidance when the surrogate is confident (in-distribution),
  \item Reduce guidance when the surrogate is uncertain (out-of-distribution early steps),
  \item Prevent adversarial hallucination caused by overconfident guidance in noisy regions.
\end{itemize}

This is directly analogous to active-learning acquisition functions, but applied to
generative trajectories.
\textbf{No prior work} has proposed uncertainty-weighted guidance in physics-constrained
diffusion for structural inverse design.

\subsection{Contribution 3: Differentiable Connectivity Discriminator}

\paragraph{What it is.}
A learned binary classifier $D_{\mathrm{conn}}(x) \in [0,1]$ predicts topological
validity.
During generation, its gradient $\nabla_{\xtilde}\log D_{\mathrm{conn}}(\xtilde)$
provides differentiable topology guidance.

\paragraph{Why it is novel.}
All prior EM inverse design work either:
\begin{itemize}
  \item Ignores connectivity entirely (most neural methods), or
  \item Applies post-hoc graph repair (which breaks differentiability).
\end{itemize}
The differentiable discriminator is analogous to discriminator guidance in GANs but
applied to topological validity rather than visual realism.

\subsection{Contribution 4: Procedurally Structured EM Dataset}

\paragraph{What it is.}
A dataset of 100k--500k physically meaningful EM structures generated from parameterised
RF primitives (not random binary images).

\paragraph{Why it is novel.}
The dataset generation methodology itself is a scientific contribution.
A random binary pixel dataset is essentially useless for training a physics-aware
generative model: it overwhelmingly samples physically degenerate configurations.
The procedural approach ensures systematic coverage of resonant, filtering, and coupling
behaviours with controlled diversity.

\subsection{Contribution 5: End-to-End Manufacturable Synthesis Pipeline}

A fully integrated pipeline:
\[
  \text{Spectral specification}
  \;\to\; \text{D3PM sampling}
  \;\to\; \text{Physics guidance}
  \;\to\; \text{Topology validation}
  \;\to\; \text{DRC check}
  \;\to\; \text{Fabrication export}
\]

No prior work handles both spectral accuracy \emph{and} manufacturability in a single
differentiable system.

\subsection{Contribution 6: Rigorous Discrete Diffusion with Continuous Physics Guidance}

\paragraph{What it is.}
A mathematically correct reconciliation of discrete diffusion (Bernoulli/categorical
state space) with continuous differentiable physics guidance via an explicit relaxation
mechanism.

\paragraph{Why it is novel.}
The v1 document contained a fundamental mathematical error—applying continuous CFG to
discrete diffusion.
The corrected formulation (Sections~\ref{sec:cfg} and~\ref{sec:guidance}) provides a
mathematically sound algorithm that is new in the discrete diffusion literature.

% ============================================================
\section{Dataset Generation Framework}
\label{sec:dataset}
% ============================================================

\subsection{Why a Random Binary Dataset Fails}

Uniform random binary layouts overwhelmingly produce:
disconnected conductors, broadband scattering noise, weak resonances, and degenerate
transmission behaviour.

Training on such distributions causes mode collapse, poor topology priors, and
physically meaningless generations.
The manifold of \emph{useful} EM structures is vanishingly thin in random binary space.

\subsection{Procedural Generation by Physical Primitives}

Every structure in the dataset is constructed from parameterised physical primitives:

\begin{table}[H]
\centering
\caption{Dataset primitive taxonomy with physical motivations.}
\label{tab:dataset_primitives}
\begin{tabularx}{\linewidth}{lXl}
\toprule
\textbf{Category} & \textbf{Primitive} & \textbf{Dominant Behaviour} \\
\midrule
Passband  & Microstrip line, tapered impedance transition & Low-loss transmission \\
Stopband  & $\lambda/4$ shunt stub, $\lambda/2$ resonator, notch & Band rejection at $f_0$ \\
Bandpass  & Coupled $\lambda/2$ resonators, interdigital, ring resonator & Selective transmission \\
Coupling  & Edge-coupled lines, broadside coupling & Directional power transfer \\
\bottomrule
\end{tabularx}
\end{table}

\subsection{Stochastic Perturbation Strategy}
\label{sec:perturbation}

\begin{lstlisting}[caption={Pseudocode for procedural dataset generation.}]
for each base primitive P:
    geom_params = perturb(P.nominal_params,
                          sigma = 0.15 * param_range)
    P_perturbed  = render_primitive(geom_params)
    P_pixel      = rasterize(P_perturbed, grid=15x15)

    if random() < p_mutate:
        P_pixel = local_mutation(P_pixel, n_flips=U(1,5))

    if random() < p_combine:
        P_pixel = combine(P_pixel, sample_primitive())
\end{lstlisting}

\paragraph{Mutation types.}
\begin{itemize}
  \item Random pixel flips in a local $3\times3$ neighbourhood (preserves macrostructure).
  \item Edge erosion/dilation (simulates fabrication tolerances).
  \item Gap insertion (creates capacitive discontinuities).
  \item Stub addition (introduces parasitic inductance/capacitance).
\end{itemize}

\subsection{Connectivity Enforcement Algorithm}

Every generated layout must pass the following four-step validation:

\begin{enumerate}
  \item \textbf{Port identification.} Identify input/output ports at fixed positions
        (left/right edges, centre row). Mark port pixels as mandatory conductors.
  \item \textbf{Connected component analysis.} BFS/DFS from port~1.
        Check reachability to port~2.
  \item \textbf{Island detection.} Identify all connected components.
        Flag components not connected to any port as floating islands.
  \item \textbf{Validation.}
        PASS if all ports are reachable and there are no floating conductors.
        FAIL $\Rightarrow$ discard and regenerate.
\end{enumerate}

\NOTE{\textbf{Why this matters.}
OpenEMS will simulate a disconnected structure and return physically nonsensical
S-parameters (infinite impedance, zero transmission).
The dataset \emph{must} exclude all disconnected structures.}

\subsection{Electromagnetic Simulation Pipeline}

\begin{table}[H]
\centering
\caption{OpenEMS simulation parameters per structure.}
\label{tab:sim_params}
\begin{tabular}{ll}
\toprule
\textbf{Parameter} & \textbf{Value} \\
\midrule
Solver         & OpenEMS (FDTD, open-source) \\
Grid resolution & $\Delta = L/H$ (matched to physical pixel size) \\
Frequency sweep & $f_1 = 0.5$~GHz, $f_N = 20$~GHz, $N_f = 100$ (linear) \\
Substrates      & $\varepsilon_r \in \{2.2,\; 3.5,\; 4.4,\; 10.2\}$ \\
Port impedance  & $Z_0 = 50\;\Omega$ \\
Time/structure  & $\sim$1--3~min (single CPU core) \\
\bottomrule
\end{tabular}
\end{table}

\subsection{Data Storage Format}

\begin{table}[H]
\centering
\caption{HDF5 dataset schema.}
\label{tab:data_schema}
\begin{tabular}{llll}
\toprule
\textbf{Field} & \textbf{Shape} & \textbf{Type} & \textbf{Description} \\
\midrule
\texttt{layout}         & $(15,15)$   & \texttt{uint8}   & Binary pixel map \\
\texttt{S11\_mag}       & $(N_f,)$    & \texttt{float32} & Return loss magnitude \\
\texttt{S21\_mag}       & $(N_f,)$    & \texttt{float32} & Insertion loss magnitude \\
\texttt{S11\_phase}     & $(N_f,)$    & \texttt{float32} & Return loss phase (rad) \\
\texttt{S21\_phase}     & $(N_f,)$    & \texttt{float32} & Insertion loss phase (rad) \\
\texttt{substrate\_id}  & scalar      & \texttt{uint8}   & Substrate type index \\
\texttt{resonance\_f}   & $(K,)$      & \texttt{float32} & Resonance frequencies \\
\texttt{Q\_factor}      & $(K,)$      & \texttt{float32} & Q-factor at resonances \\
\texttt{validity}       & scalar      & \texttt{bool}    & Passed all physical checks \\
\bottomrule
\end{tabular}
\end{table}

\textbf{Dataset splits:} Train 80\% / Validation 10\% / Test 10\% (held-out, diverse types).

\subsection{Dataset Scale Justification}

For the prior to cover the EM manifold adequately:
\begin{itemize}
  \item $10^+$ primitive types $\times$ 1000 geometric parameter combinations
        $= 10{,}000$ base structures.
  \item $\times\;10$ perturbation variants each $= 100{,}000$ minimum.
\end{itemize}

Target: \textbf{100k--500k structures}.
At 100k the model generalises within primitive families.
At 500k cross-primitive combinations are well-represented.

% ============================================================
\section{Spectral Specification Encoder}
\label{sec:encoder}
% ============================================================

\subsection{Why Frequency Correlations Cannot Be Ignored}

The spectral specification $\mathbf{y}$ is a frequency-domain sequence with strong
internal structure:
\begin{itemize}
  \item Resonances are \emph{local} in frequency (Lorentzian line shape).
  \item Passband/stopband transitions have specific roll-off characteristics.
  \item Phase is the Hilbert transform of log-magnitude for minimum-phase systems.
\end{itemize}

A flat MLP encoder ignoring these correlations will confuse:
a resonance at 5~GHz with one at 10~GHz (different layout length scales),
a sharp vs.\ broad resonance (different Q-factor),
or a 3-pole vs.\ a 5-pole filter response.
\emph{The architecture must explicitly model frequency correlations.}

\subsection{Encoder Architecture}

\paragraph{Recommended: 1D ResNet Encoder.}

\begin{lstlisting}[caption={1D ResNet spectral encoder.}]
Input : y in R^{4 x N_f}   (S11_mag, S21_mag, S11_phase, S21_phase)
  |
Conv1D(4->32, kernel=7, stride=1, pad=3) + LayerNorm + ReLU
  |
ResBlock1D x 3  (32 ch, kernel=5)
  |
Conv1D(32->64, kernel=3, stride=2)   # downsample freq axis by 2
  |
ResBlock1D x 3  (64 ch, kernel=5)
  |
Conv1D(64->128, kernel=3, stride=2)  # downsample by 2 again
  |
GlobalAveragePool -> Linear(128 -> m)
  |
Output: c_y in R^m,  m = 128 or 256
\end{lstlisting}

\paragraph{Alternative: Lightweight Transformer Encoder.}
Tokenise the spectrum into 10-frequency chunks ($N_f/10$ tokens), apply 4-head
attention with 2 layers, and use the CLS-token output as $\cy$.
Use this if ablation shows insufficient frequency correlation capture with the ResNet.

\subsection{Conditioning Mechanism}

The latent embedding $\cy \in \mathbb{R}^m$ is injected into the denoiser via
\textbf{Adaptive Layer Normalisation (AdaLN)}:

\begin{equation}
  \mathrm{AdaLN}(h,\, \cy,\, t)
  = \bigl(1 + \gamma([\cy;\, e_t])\bigr)
    \cdot \mathrm{LayerNorm}(h)
    + \beta([\cy;\, e_t])
  \label{eq:adaLN}
\end{equation}

where $e_t$ is the sinusoidal timestep embedding and $\gamma, \beta$ are learned linear
projections.
This is the DiT-style conditioning (Peebles \& Xie, 2022) and has proven stable
and efficient across many architectures.

% ============================================================
\section{Forward Surrogate Physics Model}
\label{sec:surrogate}
% ============================================================

\subsection{Role and Requirements}

The surrogate $\hat{\mathcal{F}} : \{0,1\}^{\HW} \to \mathbb{R}^d$ must satisfy four
requirements:

\begin{enumerate}
  \item \textbf{Predictive accuracy.}
        $\|\hat{\mathcal{F}}(x) - \FEM(x)\|_2 < \epsilon_{\mathrm{surr}}$ on test structures.
  \item \textbf{Gradient fidelity.}
        $\|\nabla_{\xtilde}\hat{\mathcal{F}}(\xtilde) -
           \nabla_{\xtilde}\FEM^{\mathrm{FD}}(\xtilde)\|_2 < \epsilon_{\mathrm{grad}}$.
  \item \textbf{Calibrated uncertainty.}
        Ensemble variance $\sigmahat^2(x)$ is calibrated to actual prediction error.
  \item \textbf{Fast inference.}
        ${<}\,10$~ms per evaluation on GPU.
\end{enumerate}

\subsection{Architecture}

\begin{lstlisting}[caption={Physics-aware CNN surrogate.}]
Input: x_tilde in [0,1]^{1 x 15 x 15}  (continuous relaxed layout)
  |
ConvBlock(1->32, 3x3, pad=1) + BN + ReLU
  |
ResBlock x 3  (32 ch)
  |
ConvBlock(32->64, 3x3, stride=2, pad=1)  -->  8x8
  |
ResBlock x 3  (64 ch)
  |
ConvBlock(64->128, 3x3, stride=2, pad=1)  -->  4x4
  |
ResBlock x 2  (128 ch)
  |
GlobalAveragePool -> Flatten -> Linear(128 -> 512)
  |
Linear(512 -> d)          # d = 4 x N_f
  |
Reshape -> Y_hat in R^{4 x N_f}   (S11, S21 mag + phase)
\end{lstlisting}

\NOTE{At $15\times15$ input with stride-2 downsampling, two stages produce $4\times4$
feature maps.  Do not apply a third downsampling stage at this resolution—it collapses
spatial information.}

\subsection{Ensemble Training Protocol}

Train $K = 5$ independent surrogates with different random seeds and data orderings.

\begin{align}
  \hat{\mu}(x)    &= \frac{1}{K}\sum_{k=1}^K \hat{\mathcal{F}}_k(x)
  \label{eq:ensemble_mean}\\
  \hat{\sigma}^2(x) &= \frac{1}{K-1}\sum_{k=1}^K
                       \bigl(\hat{\mathcal{F}}_k(x) - \hat{\mu}(x)\bigr)^2
  \label{eq:ensemble_var}
\end{align}

The ensemble variance $\hat{\sigma}^2$ provides calibrated uncertainty for the
guidance mechanism (Section~\ref{sec:guidance}).

\subsection{Gradient Fidelity Validation Protocol}
\label{sec:grad_fidelity}

The gradient $\nabla_{\xtilde}\hat{\mathcal{F}}(\xtilde)$ is critical.
Validate it as follows:

\begin{enumerate}
  \item Sample 1000 test layouts $x^{(i)}$.
  \item Compute surrogate gradient:
        $g_s^{(i)} = \nabla_{\xtilde}\hat{\mathcal{F}}(\xtilde^{(i)})$.
  \item Compute finite-difference EM gradient:
        $g_{\mathrm{FD},j}^{(i)} =
          \bigl(\FEM(x^{(i)} + \delta e_j) - \FEM(x^{(i)})\bigr)/\delta$.
  \item Report \textbf{cosine similarity}
        $\cos(g_s, g_{\mathrm{FD}})$
        and \textbf{gradient magnitude ratio}
        $\|g_s\|/\|g_{\mathrm{FD}}\|$.
\end{enumerate}

\textbf{Acceptance criterion:} Mean cosine similarity $> 0.7$ over the test set.
Below this threshold, guidance will be unreliable or actively harmful.

% ============================================================
\section{Generative Topology Model — Architecture and Theory}
\label{sec:generative}
% ============================================================

\subsection{Why Discrete Diffusion?}

EM layouts are \emph{binary}.
Gaussian (continuous) diffusion corrupts data by adding Gaussian noise—applied to
$\{0,1\}$ values, intermediate states become meaningless real numbers.
The correct inductive bias is a corruption process that operates natively in
$\{0, 1\}$ (or an extended ternary state space).

\begin{table}[H]
\centering
\caption{Comparison of generative frameworks for binary EM layouts.}
\label{tab:gen_compare}
\begin{tabular}{lp{5.5cm}l}
\toprule
\textbf{Framework} & \textbf{Core Mechanism} & \textbf{Suitability} \\
\midrule
D3PM (Austin et al., 2021)
  & Categorical corruption (absorbing, uniform, or Gaussian)
  & \textbf{High} \\
MDLM / Masked Diffusion (Shi et al., 2024)
  & Masking with BERT-style unmasking
  & \textbf{High} \\
Discrete Flow Matching (Campbell et al., 2024)
  & Continuous-time flows over discrete space
  & \textbf{Very High} \\
Continuous DDPM on binary
  & Gaussian corruption of $\{0,1\}$ values
  & \textbf{Low — incorrect} \\
\bottomrule
\end{tabular}
\end{table}

\textbf{Recommendation:} Use \textbf{D3PM with absorbing-state corruption} as the
primary framework.
The absorbing (MASK) state corresponds to a MASK token providing clean interpolation
between known and unknown pixels—directly interpretable as a conditional inpainting
problem.

\subsection{D3PM Forward Process (Correct Formulation)}

For binary layouts, we extend to a ternary space $\{0, 1, m\}$ where $m$ is the MASK
token.
The forward transition matrix is:

\begin{equation}
  Q_t =
  \begin{pmatrix}
    1 - \beta_t & 0           & \beta_t \\
    0           & 1 - \beta_t & \beta_t \\
    0           & 0           & 1
  \end{pmatrix}
  \label{eq:Qt}
\end{equation}

The marginal corruption at time $t$ (probability a pixel is masked):

\begin{equation}
  \bar{q}(x_t = m \mid x_0) = 1 - (1 - \bar{\beta}_t),
  \qquad
  \bar{\beta}_t = 1 - \prod_{s=1}^t (1 - \beta_s)
  \label{eq:marginal}
\end{equation}

At $t = 0$: all pixels are known.
At $t = T$: all pixels are MASK.
This gives a clean, interpretable masking schedule.

\paragraph{Reverse posterior (used in the ELBO):}

\begin{equation}
  q(x_{t-1} \mid x_t, x_0)
  = \frac{q(x_t \mid x_{t-1})\; q(x_{t-1} \mid x_0)}{q(x_t \mid x_0)}
  \label{eq:reverse_posterior}
\end{equation}

This is analytically tractable for the absorbing transition.

\subsection{D3PM Training Objective (ELBO)}

\begin{equation}
  -\mathrm{ELBO}
  = \underbrace{\mathbb{E}\bigl[-\log p_\theta(x_0 \mid x_1)\bigr]}_{\text{reconstruction}}
  + \sum_{t=2}^{T}
    \underbrace{
      \mathbb{E}\bigl[
        D_{\mathrm{KL}}\!\bigl(
          q(x_{t-1} \mid x_t, x_0)
          \;\|\;
          p_\theta(x_{t-1} \mid x_t)
        \bigr)
      \bigr]
    }_{\text{denoising objective at timestep } t}
  \label{eq:elbo}
\end{equation}

The model $p_\theta(x_{t-1} \mid x_t, \cy)$ predicts a categorical distribution over
$\{0, 1, m\}$ for each pixel, conditioned on the noisy layout and spectral embedding.

\subsection{Denoiser Architecture}

\begin{lstlisting}[caption={Topology-Aware U-Net with spectral conditioning.}]
Input: x_t in {0,1,m}^{15x15}  (learned 3D pixel embedding)
       c_y  in R^m              (spectral conditioning vector)
       t                        (timestep: sinusoidal + MLP)

Encoder:
  E0: Conv(embed->64, 3x3, pad=1) + AdaLN(c_y,t) + GELU
  E1: ResBlock(64, c_y, t) x 2
  E2: Conv(64->128, 3x3, stride=2) + AdaLN(c_y,t)  # -> 8x8
  E3: ResBlock(128, c_y, t) x 2

Bottleneck:
  B:  SelfAttention(128) + AdaLN(c_y,t)   # 8x8 = 64 tokens

Decoder:
  D3: ResBlock(128, c_y, t) x 2
  D2: TransposeConv(128->64, 2x2, stride=2, output_size=(15,15))
      + concat(E1) + AdaLN(c_y,t)
  D1: ResBlock(64, c_y, t) x 2

Output:
  Linear(64 -> 3)    # logits over {0, 1, MASK} per pixel
  Softmax -> categorical distribution
\end{lstlisting}

\NOTE{\textbf{Odd-resolution issue.}
The $15 \to 8$ stride-2 downsampling followed by upsampling back to $15\times15$
requires explicit \texttt{output\_size=(15,15)} in \texttt{TransposeConv2d}.
Always set this parameter explicitly to avoid off-by-one errors.}

\subsection{AdaLN Conditioning Details}

For each AdaLN layer with conditioning $(\cy, t)$:

\begin{equation}
  \mathrm{AdaLN}(h, \cy, t)
  = \bigl(1 + \gamma([\cy;\, e_t])\bigr)
    \cdot \mathrm{LayerNorm}(h)
    + \beta([\cy;\, e_t])
  \label{eq:adaLN2}
\end{equation}

where $[\cy;\, e_t] \in \mathbb{R}^{m + d_t}$ is the concatenated condition and
$\gamma, \beta$ are learned linear projections.

% ============================================================
\section{Classifier-Free Guidance in Discrete Space}
\label{sec:cfg}
% ============================================================

\subsection{The V1 Error}

\WARN{
Version~1 copied the continuous DDPM CFG formula verbatim:
\[
  \hat{\epsilon} = \epsilon_\theta(x_t, t, \emptyset)
    + w\bigl(\epsilon_\theta(x_t, t, \cy) - \epsilon_\theta(x_t, t, \emptyset)\bigr)
\]
\textbf{This is mathematically wrong for discrete diffusion.}
In continuous DDPM, $\epsilon_\theta$ predicts a Gaussian noise vector and guidance
modifies the Gaussian mean.
In discrete diffusion there is no noise vector—the model predicts a
\emph{categorical distribution} over pixel states.
Modifying a noise vector in a categorical distribution is undefined.
}

\subsection{Correct CFG for Discrete Diffusion}

The correct formulation operates at the \textbf{log-probability level}.

Let $p_\theta(x_0 \mid x_t, \cy)$ and $p_\theta(x_0 \mid x_t, \emptyset)$ be the
conditional and unconditional denoiser distributions respectively.
The guided distribution is:

\begin{equation}
  \log \tilde{p}_\theta(x_0 \mid x_t, \cy)
  = (1 + w)\log p_\theta(x_0 \mid x_t, \cy)
  - w\log p_\theta(x_0 \mid x_t, \emptyset)
  \label{eq:disc_cfg}
\end{equation}

normalised to a valid probability.
Since pixels are predicted independently (per-pixel categorical), this factorises as:

\begin{equation}
  \tilde{p}_\theta(x_0^{(ij)} \mid x_t, \cy)
  \propto
  p_\theta(x_0^{(ij)} \mid x_t, \cy)^{1+w}
  \cdot
  p_\theta(x_0^{(ij)} \mid x_t, \emptyset)^{-w}
  \label{eq:disc_cfg_factored}
\end{equation}

\FIX{
Training uses condition dropout probability $p_{\mathrm{drop}} = 0.1$--$0.2$:
during training, the spectral condition $\cy$ is randomly replaced by a null embedding
$\emptyset$.
At inference, both the conditional and unconditional forward passes are run and combined
via~\eqref{eq:disc_cfg_factored}.
}

\subsection{Guidance Scale Selection}

\begin{table}[H]
\centering
\caption{Effect of CFG guidance weight $w$.}
\label{tab:cfg_w}
\begin{tabular}{cll}
\toprule
$w$ & Behaviour & Use Case \\
\midrule
$0$   & Unconditional sampling (max diversity)      & Exploration \\
$1$--$3$ & Balanced (recommended starting point)   & General synthesis \\
$2$--$4$ & High spectral fidelity (expected optimal for RF) & Precise specs \\
$>5$  & Strong guidance (risk of mode collapse)     & Avoid unless necessary \\
\bottomrule
\end{tabular}
\end{table}

% ============================================================
\section{Physics-Guided Sampling}
\label{sec:guidance}
% ============================================================

\subsection{The V1 Error}

\WARN{
Version~1 proposed:
\[
  \hat{\epsilon} = \tilde{\epsilon}
    - \alpha_t \nabla_{x_t} \mathcal{L}_{\mathrm{physics}}
\]
\textbf{This is undefined.}
$x_t \in \{0, 1, m\}$ is a discrete/ternary token.
Gradients of continuous functions w.r.t.\ discrete tokens are not defined.
}

\subsection{Corrected: Guidance Through Predicted $\hat{x}_0$}

\paragraph{Key insight.}
The denoiser predicts the clean layout $\hat{x}_0 = f_\theta(x_t, t, \cy)$.
If we take the \emph{expected value} under the predicted distribution, we get a
continuous quantity:

\begin{equation}
  \hat{x}_0 = \mathbb{E}_{p_\theta(x_0 \mid x_t, \cy)}[x_0]
  = \sum_{v \in \{0,1,m\}} v \cdot p_\theta(x_0 = v \mid x_t, \cy)
  \label{eq:x0hat}
\end{equation}

This gives $\hat{x}_0 \in [0,1]^{\HW}$ (the per-pixel probability of being conductor),
which is \emph{differentiable} with respect to the denoiser parameters and logits.

\paragraph{Physics loss on predicted $\hat{x}_0$:}

\begin{equation}
  \mathcal{L}_{\mathrm{physics}}
  = \left\|\hat{\mathcal{F}}(\hat{x}_0) - \ystar\right\|_2^2
  \label{eq:lphysics}
\end{equation}

\paragraph{Gradient-guided update to the logits:}

Let $\ell_t \in \mathbb{R}^{H \times W \times 3}$ be the raw denoiser logits.
The physics guidance modifies these logits, not $x_t$:

\begin{equation}
  \tilde{\ell}_t
  = \ell_t
  - \alpha_t \nabla_{\ell_t}
    \mathcal{L}_{\mathrm{physics}}\!\left(\hat{x}_0(\ell_t),\, \ystar\right)
  \label{eq:logit_guidance}
\end{equation}

where $\hat{x}_0(\ell_t) = \sigma(\ell_t^{[1]})$ (probability of pixel $= 1$ from
softmax of logits).
This is differentiable end-to-end:
$\ell_t \;\to\; \hat{x}_0 \;\to\; \hat{\mathcal{F}}(\hat{x}_0) \;\to\;
\mathcal{L}_{\mathrm{physics}}$.

\subsection{Uncertainty-Weighted Guidance}

\begin{equation}
  \alpha_t
  = \frac{\alphamax}{\sigmahat(\hat{x}_0) + \epsilon}
    \cdot \eta_t,
  \qquad
  \eta_t = \left(1 - \frac{t}{T}\right)^2
  \label{eq:alpha_t}
\end{equation}

where:
\begin{itemize}
  \item $\sigmahat(\hat{x}_0) = \sqrt{\mathrm{mean}_f(\hat{\sigma}^2_{\mathrm{ensemble}}(\hat{x}_0))}$
        is the mean ensemble uncertainty across frequencies,
  \item $\eta_t = (1 - t/T)^2$ is a timestep-decay schedule (guidance activates as
        denoising completes),
  \item $\alphamax$ is a hyperparameter controlling maximum step size, and
  \item $\epsilon = 0.01$ prevents division by zero.
\end{itemize}

This formula guarantees:
\begin{enumerate}
  \item No guidance at early timesteps (noisy layouts → meaningless surrogate evaluations).
  \item Strong guidance at late timesteps (topology is nearly crystallised → accurate surrogate).
  \item Automatic step reduction when the surrogate is uncertain (prevents adversarial guidance).
\end{enumerate}

\subsection{Late-Stage Activation Threshold}

Guidance activates only when $t < \tthresh$.
Expected optimal: $\tthresh = 0.3T$--$0.5T$.

\paragraph{Rationale.}
At early timesteps ($t > \tthresh$), most pixels are MASK tokens.
The predicted $\hat{x}_0$ is highly uncertain and the surrogate gradient points in a
physically meaningless direction.
At late timesteps, the topology is largely decided and small corrections for spectral
accuracy are effective and reliable.

\subsection{Complete Physics-Guided Sampling Algorithm}

\begin{algorithm}[H]
\caption{Physics-Guided D3PM Reverse Sampling}
\label{alg:sampling}
\KwIn{Target spec $\ystar$, spectral embedding $\cy$, trained denoiser $f_\theta$,
      trained surrogate ensemble $\hat{\mathcal{F}}$, guidance weight $w$,
      threshold $\tthresh$}
\KwOut{Generated binary layout $x_0$}

$x_T \leftarrow$ all MASK tokens\;

\For{$t = T, T-1, \ldots, 1$}{
  Compute logits: $\ell_t \leftarrow f_\theta(x_t,\, t,\, \cy)$\;
  Compute $\hat{x}_0 \leftarrow \sigma(\ell_t^{[1]})$ \tcp*{predicted layout}

  \If{$t < \tthresh$}{
    $g \leftarrow \nabla_{\ell_t}\mathcal{L}_{\mathrm{guided}}(\hat{x}_0)$\;
    $\tilde{\ell}_t \leftarrow \ell_t - \alpha_t\, g$ \tcp*{physics + topo guidance}
  }\Else{
    $\tilde{\ell}_t \leftarrow \ell_t$\;
  }

  Apply CFG via eq.~\eqref{eq:disc_cfg_factored}\;
  Sample: $x_{t-1} \sim p_\theta(x_{t-1} \mid x_t,\, \tilde{\ell}_t)$\;
}
\Return $x_0$\;
\end{algorithm}

% ============================================================
\section{Differentiable Topology Constraints}
\label{sec:topology}
% ============================================================

\subsection{Why Post-Hoc Repair Fails}

Every prior neural EM inverse design method either ignores connectivity entirely or
applies graph repair after generation.
Post-hoc repair \emph{breaks differentiability}: the repair step cannot be
backpropagated through, so the generative model never learns to avoid disconnected
topologies in the first place.
The fix is to embed topology validity as a differentiable objective \emph{inside} the
generation loop.

\subsection{Connectivity Discriminator — Training}

\paragraph{Supervision signal.}
For each layout $x$ in the dataset, run DFS to determine:
\begin{itemize}
  \item $y_{\mathrm{conn}} = 1$: Layout is connected (port~1 $\to$ port~2 path exists,
        no isolated conductor islands).
  \item $y_{\mathrm{conn}} = 0$: Layout is disconnected or has no inter-port path.
\end{itemize}

\paragraph{Architecture.}
\begin{lstlisting}[caption={Connectivity discriminator architecture.}]
Input: x in [0,1]^{15x15}
  |
Conv(1->32, 3x3) -> BN -> ReLU
  |
ResBlock x 3  (32 ch)
  |
GlobalAveragePool
  |
Linear(32->16) -> ReLU -> Linear(16->1) -> Sigmoid
  |
Output: D_conn(x) in [0,1]
\end{lstlisting}

\paragraph{Loss.}
Binary cross-entropy with balanced class weights.
Disconnected structures are rare in the procedural dataset, so explicitly oversample
them when training the discriminator.

\subsection{Connectivity Guidance During Sampling}

Add a topology loss to the guided objective:

\begin{equation}
  \mathcal{L}_{\mathrm{topo}} = -\log D_{\mathrm{conn}}(\hat{x}_0)
  \label{eq:ltopo}
\end{equation}

The combined guidance loss at each step:

\begin{equation}
  \mathcal{L}_{\mathrm{guided}}
  = \mathcal{L}_{\mathrm{physics}}
  + \lambda_{\mathrm{topo}}\, \mathcal{L}_{\mathrm{topo}}
  + \lambda_{\mathrm{mfg}}\, \mathcal{L}_{\mathrm{mfg}}
  \label{eq:lguided}
\end{equation}

Apply a ramp-up schedule on $\lambda_{\mathrm{topo}}$: weak at early timesteps
(layout is mostly MASK), strong at late timesteps (topology nearly decided).

\subsection{Port Connectivity Encoding}

Port positions are fixed by design (e.g., left-edge centre pixel = port~1,
right-edge centre pixel = port~2).
Encode port positions as a second binary channel appended to the layout:

\begin{equation}
  x_{\mathrm{aug}} = [x,\, p] \in \{0,1,\textsc{mask}\}^{2 \times 15 \times 15}
\end{equation}

where $p \in \{0,1\}^{15 \times 15}$ has $1$s only at port pixels.
This informs all components—surrogate, discriminator, denoiser—of the port structure.

% ============================================================
\section{Differentiable Binarisation}
\label{sec:binarisation}
% ============================================================

\subsection{The Discretisation Problem}

Hard thresholding ($x = \mathbbm{1}[\xtilde > 0.5]$) is not differentiable—no gradient
flows back through it.
Yet the surrogate and connectivity discriminator both require a real-valued relaxation
for gradient-based guidance.

\subsection{Annealed Gumbel-Sigmoid Binarisation}

Use the Gumbel-Sigmoid trick (Maddison et al., 2016; Jang et al., 2016) with
temperature annealing:

\begin{equation}
  \xtilde_{ij}
  = \sigma\!\left(
      \frac{\log p_{ij} - \log(1 - p_{ij}) + g_{ij}}{\tau_t}
    \right)
  \label{eq:gumbel}
\end{equation}

where:
\begin{itemize}
  \item $p_{ij} \in [0,1]$ is the predicted probability of pixel $(i,j)$ being conductor,
  \item $g_{ij} \sim \mathrm{Gumbel}(0,1)$ provides reparameterisation noise, and
  \item $\tau_t > 0$ is the temperature, annealed:
        $\tau_t = \tau_{\max} \cdot (\tau_{\min}/\tau_{\max})^{t/T}$.
\end{itemize}

\paragraph{At inference.}
$\tau \to 0 \;\Rightarrow\; \xtilde_{ij} \to \mathbbm{1}[p_{ij} > 0.5]$
(hard binarisation).

\paragraph{During training.}
$\tau = \tau_{\mathrm{initial}} = 1.0$, so soft gradients flow through.

\paragraph{Straight-Through Estimator (STE) as fallback.}
Forward pass uses hard binarisation;
backward pass passes gradients as if $x_{ij} = \xtilde_{ij}$.
Less principled but highly effective for binary variables in practice.

\NOTE{\textbf{Recommendation.}
Use Gumbel-Sigmoid with temperature annealing for the main surrogate training.
Use STE for gradient guidance during inference (simpler, no Gumbel noise to tune).}

% ============================================================
\section{Manufacturability Constraints}
\label{sec:mfg}
% ============================================================

\subsection{DRC Rules for Binary Pixel Layouts}

At the physical scale of the $15\times15$ grid (pixel size $\Delta$~mm):

\begin{table}[H]
\centering
\caption{Manufacturing design rules and their enforcement strategies.}
\label{tab:drc}
\begin{tabular}{lll}
\toprule
\textbf{Constraint} & \textbf{Definition} & \textbf{Enforcement} \\
\midrule
Minimum trace width  & $w_{\min} = 2$ pixels        & Penalise isolated single-pixel conductors \\
Minimum spacing      & $s_{\min} = 2$ pixels        & Penalise conductors $<2$ pixels apart \\
No floating islands  & No isolated conductor clusters & Connectivity discriminator \\
No acute corners     & No $1\times1$ corner conductors & Morphological erosion check \\
\bottomrule
\end{tabular}
\end{table}

\subsection{Differentiable DRC Loss}

\begin{equation}
  \mathcal{L}_{\mathrm{mfg}}
  = \lambda_1\,\mathcal{L}_{\mathrm{width}}
  + \lambda_2\,\mathcal{L}_{\mathrm{spacing}}
  + \lambda_3\,\mathcal{L}_{\mathrm{islands}}
  \label{eq:lmfg}
\end{equation}

\paragraph{Minimum width loss} (penalise thin traces):
\begin{equation}
  \mathcal{L}_{\mathrm{width}}
  = \sum_{(i,j)} \xtilde_{ij}
    \cdot \max\!\left(0,\; 1 - \sum_{(i',j') \in \mathcal{N}(i,j)} \xtilde_{i'j'}\right)
  \label{eq:lwidth}
\end{equation}

where $\mathcal{N}(i,j)$ is the 4-connected neighbourhood.

\paragraph{Minimum spacing loss} (penalise nearby isolated conductors):
\begin{equation}
  \mathcal{L}_{\mathrm{spacing}}
  = \sum_{(i,j)} \xtilde_{ij}
    \cdot (1 - \xtilde_{\mathrm{eroded},(i,j)})
    \cdot D_{\mathrm{nearest}}(\xtilde,\, i,\, j)
  \label{eq:lspacing}
\end{equation}

where $\xtilde_{\mathrm{eroded}}$ is a differentiable morphological erosion.

\paragraph{Islands loss.}
Directly reuse the connectivity discriminator:
$\mathcal{L}_{\mathrm{islands}} = -\log D_{\mathrm{conn}}(\hat{x}_0)$.

% ============================================================
\section{Full Training and Inference Objectives}
\label{sec:training}
% ============================================================

\subsection{Surrogate Training Objective}

\begin{equation}
  \mathcal{L}_{\mathrm{surrogate}}
  = \underbrace{\|\yhat - \mathbf{y}\|_2^2}_{\mathrm{MSE}}
  + \lambda_{\mathrm{pass}} \underbrace{\mathcal{L}_{\mathrm{pass}}}_{\mathrm{passivity}}
  + \lambda_{\mathrm{recip}} \underbrace{\mathcal{L}_{\mathrm{recip}}}_{\mathrm{reciprocity}}
  + \lambda_{KK} \underbrace{\mathcal{L}_{KK}}_{\mathrm{causality}}
  + \lambda_{\mathrm{smooth}} \underbrace{\mathcal{L}_{\mathrm{smooth}}}_{\mathrm{smoothness}}
  \label{eq:lsurr2}
\end{equation}

\subsection{Denoiser Training Objective}

\begin{equation}
  \mathcal{L}_{\mathrm{denoiser}}
  = \underbrace{-\mathrm{ELBO}}_{\text{D3PM ELBO}}
  + \lambda_{\mathrm{aux}}
    \underbrace{\mathcal{L}_{x_0\,\mathrm{pred}}}_{\text{auxiliary } x_0 \text{ prediction}}
  \label{eq:ldenoiser}
\end{equation}

The auxiliary $x_0$ prediction loss directly penalises error in the expected-$\hat{x}_0$
output, providing a clean gradient signal for the surrogate guidance pathway.

\subsection{Inference Algorithm (Complete)}

At each denoising step $t$ from $T$ to $1$:
\begin{enumerate}
  \item Compute denoiser output: $\ell_t = f_\theta(x_t, t, \cy)$.
  \item Compute $\hat{x}_0 = \mathrm{softmax}(\ell_t)[:,:,1]$
        (probability of pixel $= 1$).
  \item If $t < \tthresh$: compute guidance gradient
        $g = \nabla_{\ell_t}\mathcal{L}_{\mathrm{guided}}(\hat{x}_0)$.
  \item Apply guided logits: $\tilde{\ell}_t = \ell_t - \alpha_t\, g$.
  \item Sample $x_{t-1} \sim p_\theta(x_{t-1} \mid x_t, \tilde{\ell}_t)$
        using the reverse posterior.
\end{enumerate}

\subsection{Loss Weight Hyperparameters}

\begin{table}[H]
\centering
\caption{Starting loss weights (to be tuned via ablation).}
\label{tab:lossweights}
\begin{tabular}{lcl}
\toprule
\textbf{Weight} & \textbf{Value} & \textbf{Scope} \\
\midrule
$\lambda_{\mathrm{pass}}$   & 0.10 & Surrogate \\
$\lambda_{\mathrm{recip}}$  & 0.05 & Surrogate \\
$\lambda_{KK}$              & 0.05 & Surrogate \\
$\lambda_{\mathrm{smooth}}$ & 0.02 & Surrogate \\
$\lambda_{\mathrm{topo}}$   & 1.00 & Guided sampling \\
$\lambda_{\mathrm{mfg}}$    & 0.50 & Guided sampling \\
$\alphamax$                 & 0.10 & Guidance step size \\
$w$ (CFG)                   & 2.00 & Inference scale \\
\bottomrule
\end{tabular}
\end{table}

% ============================================================
\section{Vulnerability Analysis and Mitigations}
\label{sec:vulnerabilities}
% ============================================================

\subsection{Critical Vulnerabilities — V1 Gaps Now Resolved}

\begin{tcolorbox}[fixbox, title=V-CRIT-1: CFG Formulation for Discrete Diffusion]
\textbf{V1 state:} Continuous CFG formula applied verbatim to discrete diffusion.\\
\textbf{Problem:} Mathematically incorrect; modifying a noise vector in discrete space
is undefined.\\
\textbf{V2 fix:} Log-probability-domain CFG (Section~\ref{sec:cfg}).\\
\textbf{Status:} \textcolor{okgreen}{RESOLVED}
\end{tcolorbox}

\begin{tcolorbox}[fixbox, title=V-CRIT-2: Gradient Guidance on Discrete Variables]
\textbf{V1 state:}
$\hat{\epsilon} = \tilde{\epsilon} - \alpha_t \nabla_{x_t}\mathcal{L}_{\mathrm{physics}}$
applied to discrete $x_t$.\\
\textbf{Problem:} $x_t \in \{0,1,m\}$ — gradients undefined.\\
\textbf{V2 fix:} Guidance applied to denoiser logits via predicted $\hat{x}_0$
(Section~\ref{sec:guidance}).\\
\textbf{Status:} \textcolor{okgreen}{RESOLVED}
\end{tcolorbox}

\begin{tcolorbox}[fixbox, title=V-CRIT-3: Missing Kramers--Kronig Causality]
\textbf{V1 state:} No causality constraint in surrogate training.\\
\textbf{Problem:} Surrogate can learn physically impossible non-causal responses.\\
\textbf{V2 fix:} KK loss added to surrogate objective (Section~\ref{sec:physics}).\\
\textbf{Status:} \textcolor{okgreen}{RESOLVED}
\end{tcolorbox}

\begin{tcolorbox}[fixbox, title=V-CRIT-4: Incomplete Passivity Constraint]
\textbf{V1 state:} $|S_{11}|^2 + |S_{21}|^2 \leq 1$ (2-port scalar only).\\
\textbf{Problem:} Fails for multi-port extension; does not use matrix form.\\
\textbf{V2 fix:} Full matrix passivity constraint via maximum eigenvalue
(Section~\ref{sec:smatrix}).\\
\textbf{Status:} \textcolor{okgreen}{RESOLVED}
\end{tcolorbox}

\subsection{Major Vulnerabilities — Active Mitigations Required}

\paragraph{V-MAJ-1: Resolution Limitation ($15\times15$).}
\textit{Problem:} Real structures often need $32\times32$ to $128\times128$.\\
\textit{Mitigation:}
The frequency--resolution co-design table (Section~\ref{sec:resfcodesign}) shows that
$15\times15$ is valid for 2--18~GHz on 5--10~mm substrates—covering cellular, WiFi,
and satellite bands.
Include a $32\times32$ scaling experiment in the evaluation.

\paragraph{V-MAJ-2: Surrogate Gradient Fidelity.}
\textit{Problem:} CNN surrogate gradients near topology discontinuities may be noisy.\\
\textit{Mitigation:}
\begin{itemize}
  \item Surrogate trained on continuous relaxed inputs (not hard binary).
  \item Explicit gradient fidelity validation (Section~\ref{sec:grad_fidelity}).
  \item Gradient clipping: $\|g\|_2 \leq g_{\max} = 1.0$.
  \item Uncertainty-weighted guidance automatically reduces step size when uncertain.
\end{itemize}

\paragraph{V-MAJ-3: Dataset Simulation Compute.}
\textit{Problem:} 100k--500k OpenEMS simulations $\times$ 1--3~min = potentially months
on a single machine.\\
\textit{Mitigation:} Parallelise across 50--100 CPU cores (cloud HPC).
At 50 cores: $\sim$5.6~days for 200k structures.
See Section~\ref{sec:compute} for full compute budget.

\paragraph{V-MAJ-4: Dataset Distributional Bias.}
\textit{Problem:} Procedural generation may not cover all relevant topologies.\\
\textit{Mitigation:}
Include 10\% ``challenge structures'' from random-binary-plus-repair.
Evaluate on held-out primitive types not seen during training.
Report structure uniqueness and spectral diversity metrics.

\paragraph{V-MAJ-5: Multi-Objective Loss Instability.}
\textit{Problem:} Fixed $\lambda$ weights may not balance well across training phases.\\
\textit{Mitigation:}
Use GradNorm (Chen et al., 2018) to normalise gradient magnitudes across losses,
or a simple schedule that gradually increases DRC constraint weight.

\subsection{Minor Vulnerabilities}

\begin{itemize}
  \item \textbf{V-MIN-1: Missing neural baselines.}
        Add conditional VAE and conditional GAN baselines (now in evaluation plan,
        Section~\ref{sec:eval}).
  \item \textbf{V-MIN-2: No multi-port extension.}
        Explicitly scope this work as 2-port in AAAI submission;
        include multi-port extension as future work with a clear mathematical path.
  \item \textbf{V-MIN-3: No formal diversity metric.}
        Report mean pairwise Hamming distance between generated layouts for the same
        $\ystar$, and a Wasserstein-based feature diversity score.
\end{itemize}

% ============================================================
\section{Evaluation Protocol}
\label{sec:eval}
% ============================================================

\subsection{Baselines}

\begin{table}[H]
\centering
\caption{Baselines compared against PIXEL.}
\label{tab:baselines}
\begin{tabular}{lll}
\toprule
\textbf{Baseline} & \textbf{Description} & \textbf{Reference} \\
\midrule
Det-CNN       & Deterministic CNN inverse regressor          & Standard approach \\
Det-CNN + GA  & CNN initialisation + genetic algorithm       & Common EDA pipeline \\
cVAE          & Conditional VAE with spectral encoder        & Added in V2 \\
cGAN          & Conditional GAN with spectral discriminator  & Added in V2 \\
BO            & Bayesian optimisation over layout space      & Frazier (2018) \\
Diff-TO       & Differentiable topology optimisation         & Sigmund et al.\ (2003) \\
\textbf{PIXEL}& Physics-guided discrete diffusion            & \textbf{This work} \\
\bottomrule
\end{tabular}
\end{table}

\subsection{Primary Metrics: Spectral Accuracy}

All metrics are computed after \textbf{full-wave EM verification}
(OpenEMS simulation of the generated layout)—not from surrogate predictions.

\begin{table}[H]
\centering
\caption{Spectral accuracy metrics (evaluated on 1000 test specifications).}
\label{tab:metrics_spec}
\begin{tabularx}{\linewidth}{lX}
\toprule
\textbf{Metric} & \textbf{Definition} \\
\midrule
$S_{21}$ MSE & $\frac{1}{N_f}\sum_f (|S_{21}^{\mathrm{gen}}| - |S_{21}^*|)^2$ \\
$S_{11}$ MSE & $\frac{1}{N_f}\sum_f (|S_{11}^{\mathrm{gen}}| - |S_{11}^*|)^2$ \\
Resonance freq.\ error & $|f_r^{\mathrm{gen}} - f_r^*|$ for identified resonances \\
Passband error  & MSE within specified passband \\
Stopband error  & Deviation from target $S_{21}$ in stopband \\
\bottomrule
\end{tabularx}
\end{table}

\subsection{Secondary Metrics: Topological Validity}

\begin{table}[H]
\centering
\caption{Topology and manufacturability metrics.}
\label{tab:metrics_topo}
\begin{tabular}{lll}
\toprule
\textbf{Metric} & \textbf{Definition} & \textbf{Target} \\
\midrule
Connectivity yield   & Fraction of generated layouts that are connected  & $>0.95$ \\
Island-free rate     & Fraction with no floating conductors              & $>0.98$ \\
DRC pass rate        & Fraction satisfying all manufacturing constraints & $>0.90$ \\
\bottomrule
\end{tabular}
\end{table}

\subsection{Diversity Metrics}

\begin{table}[H]
\centering
\caption{Diversity metrics.}
\label{tab:metrics_div}
\begin{tabularx}{\linewidth}{lX}
\toprule
\textbf{Metric} & \textbf{Definition} \\
\midrule
Intra-spec diversity
  & Mean pairwise pixel-wise Hamming distance between samples generated for the same $\ystar$ \\
Inter-spec coverage
  & Range of distinct layout topologies generated across the full $\ystar$ test set \\
\bottomrule
\end{tabularx}
\end{table}

\subsection{Efficiency Metrics}

\begin{table}[H]
\centering
\caption{Efficiency metrics.}
\label{tab:metrics_eff}
\begin{tabularx}{\linewidth}{lX}
\toprule
\textbf{Metric} & \textbf{Definition} \\
\midrule
Inference time           & Wall-clock time per generated layout (GPU) \\
EM verification calls    & Number of full-wave simulations per final design \\
Dataset efficiency       & Accuracy as a function of training set size \\
\bottomrule
\end{tabularx}
\end{table}

\subsection{Ablation Study Design}

\begin{table}[H]
\centering
\caption{Ablation study: each row removes or replaces one component.}
\label{tab:ablation}
\begin{tabular}{lll}
\toprule
\textbf{Ablation} & \textbf{What Is Removed/Changed} & \textbf{Tests} \\
\midrule
No physics guidance        & $\alpha_t = 0$                    & Effect of physics guidance \\
No uncertainty weighting   & $\sigmahat = \mathrm{const}$      & Effect of calibrated guidance \\
No topology constraint     & $\lambda_{\mathrm{topo}} = 0$    & Effect of connectivity discriminator \\
No DRC loss                & $\lambda_{\mathrm{mfg}} = 0$     & Effect of manufacturability \\
Continuous diffusion       & Replace D3PM with DDPM            & Discrete vs.\ continuous \\
No KK loss                 & $\lambda_{KK} = 0$               & Effect of causality constraint \\
Random dataset             & Replace procedural with random binary & Dataset quality \\
Single surrogate           & Replace ensemble with single model    & Uncertainty estimation \\
\bottomrule
\end{tabular}
\end{table}

% ============================================================
\section{Related Work and Differentiation}
\label{sec:related}
% ============================================================

\subsection{Inverse EM Design}

\begin{table}[H]
\centering
\caption{Prior work on inverse EM / nanophotonic design.}
\label{tab:related_em}
\begin{tabular}{llp{5cm}}
\toprule
\textbf{Work} & \textbf{Method} & \textbf{Limitation vs.\ PIXEL} \\
\midrule
Liu et al.\ (2018)       & Tandem network             & Deterministic; no topology handling \\
Jiang et al.\ (2019)     & cGAN for nanophotonics     & Continuous image; no physics constraints \\
So et al.\ (2020)        & DNN for metasurfaces       & Single output; no diversity \\
Unni et al.\ (2021)      & VAE for photonic design    & Latent space not physics-constrained \\
Kudyshev et al.\ (2021)  & RL for EM design           & No topology or manufacturability \\
\bottomrule
\end{tabular}
\end{table}

\subsection{Discrete Diffusion}

\begin{table}[H]
\centering
\caption{Key discrete diffusion works and their relation to PIXEL.}
\label{tab:related_diff}
\begin{tabular}{lp{4.5cm}p{4.5cm}}
\toprule
\textbf{Work} & \textbf{Contribution} & \textbf{Relation to PIXEL} \\
\midrule
D3PM (Austin et al., 2021)
  & Discrete denoising diffusion
  & PIXEL uses D3PM as backbone \\
MDLM (Shi et al., 2024)
  & Masked diffusion for language
  & Alternative backbone \\
Discrete Flow Matching (Campbell et al., 2024)
  & Flow matching over discrete spaces
  & Future upgrade path \\
GDSS (Jo et al., 2022)
  & Graph discrete diffusion for molecules
  & Conceptually related; different domain \\
\bottomrule
\end{tabular}
\end{table}

\subsection{Physics-Guided Generative Models}

\begin{table}[H]
\centering
\caption{Physics-guided generative models and their relation to PIXEL.}
\label{tab:related_phys}
\begin{tabular}{lp{4.5cm}p{4.5cm}}
\toprule
\textbf{Work} & \textbf{Contribution} & \textbf{Relation to PIXEL} \\
\midrule
DPS (Chung et al., 2022)
  & Diffusion posterior sampling for inverse problems
  & PIXEL extends this to discrete spaces \\
DiffSBDD (Schneuing et al., 2022)
  & Structure-based drug design with diffusion
  & Analogous physics guidance; different domain \\
DiT (Peebles \& Xie, 2022)
  & Diffusion Transformer with AdaLN conditioning
  & PIXEL adopts AdaLN conditioning \\
Dhariwal \& Nichol (2021)
  & Classifier guidance in continuous diffusion
  & PIXEL correctly adapts to discrete space \\
\bottomrule
\end{tabular}
\end{table}

\subsection{PIXEL's Unique Position}

The unique combination that \emph{no prior work achieves}:
\begin{enumerate}
  \item Discrete diffusion (not continuous) over binary structural layouts.
  \item Physics guidance via differentiable EM surrogate (not just visual/chemical properties).
  \item Topology validity as a differentiable constraint inside the generation loop.
  \item Uncertainty-weighted guidance step size (not fixed or timestep-only schedule).
  \item Fabrication DRC integrated into the generation objective.
\end{enumerate}

% ============================================================
\section{AAAI-2027 Positioning Strategy}
\label{sec:aaai}
% ============================================================

\subsection{Fit Analysis}

AAAI-2027 is appropriate because:
\begin{enumerate}
  \item \textbf{AI methodology.}
        The paper contributes to \emph{discrete diffusion for combinatorial inverse
        problems}—this is a methodological advance independent of the EM application.
  \item \textbf{Scientific AI.}
        AAAI has a strong record of publishing physics-informed AI
        (materials discovery, drug design, protein folding).
        EM inverse design fits squarely in this paradigm.
  \item \textbf{Novelty and completeness.}
        The paper can be fully self-contained with theory, method, dataset, and evaluation.
\end{enumerate}

\subsection{Recommended Paper Framing}

\begin{tcolorbox}[notebox, title=Primary AAAI Abstract Frame]
\textit{We present PIXEL: a physics-guided discrete diffusion framework for constrained
inverse design of electromagnetic structures.
PIXEL learns a probabilistic prior over manufacturable RF topologies and guides reverse
sampling with differentiable EM physics via an uncertainty-weighted ensemble surrogate.}
\end{tcolorbox}

\textbf{Avoid:} Framing as an EDA/circuit design paper.
AAAI reviewers are AI researchers; the EM design is the application, not the
contribution.

\paragraph{Message architecture.}
\begin{enumerate}
  \item Problem: Inverse EM design is ill-posed $\to$ probabilistic formulation required.
  \item Method: D3PM + physics guidance + topology constraints.
  \item Technical novelty: Corrected discrete CFG + uncertainty-weighted guidance +
        differentiable topology.
  \item Empirical: Outperforms all deterministic baselines across all metrics.
  \item Impact: Enables direct fabrication of AI-designed RF structures.
\end{enumerate}

\subsection{Anticipated Reviewer Concerns and Responses}

\begin{table}[H]
\centering
\caption{Anticipated reviewer concerns and prepared responses.}
\label{tab:reviewer}
\begin{tabularx}{\linewidth}{lX}
\toprule
\textbf{Concern} & \textbf{Response} \\
\midrule
``$15\times15$ is too low resolution''
  & Co-design table (Sec.~\ref{sec:resfcodesign}) shows coverage of
    2--18~GHz on 5--10~mm substrates; include $32\times32$ scaling experiment. \\
``Why not use an EM optimiser?''
  & EM optimiser requires $1000\times$ more forward evaluations per design;
    present inference-time efficiency comparison. \\
``Is the surrogate accurate enough?''
  & Report gradient fidelity validation (Sec.~\ref{sec:grad_fidelity})
    and ensemble uncertainty calibration. \\
``Dataset is procedurally generated---biased?''
  & Evaluate on held-out structure types; show generalisation beyond
    training primitives; include 10\% random-plus-repair structures. \\
``Why AAAI and not IEEE TMTT?''
  & AI methodology is the primary contribution; submit IEEE TMTT in parallel for
    the engineering community. \\
\bottomrule
\end{tabularx}
\end{table}

\subsection{Concurrent Submission Strategy}

\begin{itemize}
  \item \textbf{AAAI-2027:} AI methodology focus (discrete diffusion for constrained
        inverse design).
  \item \textbf{IEEE TMTT:} Engineering focus (RF/IC design automation, detailed
        fabrication results).
  \item \textbf{ICLR-2027 Workshop:} Physics-informed ML workshop fallback.
\end{itemize}

% ============================================================
\section{Implementation Roadmap}
\label{sec:roadmap}
% ============================================================

\subsection{Phase 1 — Dataset Generation (Weeks 1--6)}

\begin{itemize}
  \item[$\square$] Implement parametric primitive generators
        (microstrip, stubs, resonators, coupled lines, SRR).
  \item[$\square$] Implement stochastic perturbation engine.
  \item[$\square$] Implement connectivity validator (BFS-based).
  \item[$\square$] Set up OpenEMS batch simulation pipeline
        with multiprocessing parallelisation.
  \item[$\square$] Generate 50k structures as initial dataset; scale to 200k as
        compute allows.
  \item[$\square$] Validate dataset: S-parameter histograms, connectivity yield,
        physical diversity metrics.
\end{itemize}

\subsection{Phase 2 — Surrogate Model (Weeks 5--9, overlaps Phase 1)}

\begin{itemize}
  \item[$\square$] Implement CNN surrogate architecture (Section~\ref{sec:surrogate}).
  \item[$\square$] Train ensemble of $K=5$ surrogates.
  \item[$\square$] Validate $S$-parameter MSE on test set.
  \item[$\square$] Validate gradient fidelity (cosine similarity $> 0.7$).
  \item[$\square$] Add KK, passivity, and reciprocity losses.
  \item[$\square$] Profile inference latency (target: ${<}\,10$~ms on GPU).
\end{itemize}

\subsection{Phase 3 — Denoiser Model (Weeks 8--14)}

\begin{itemize}
  \item[$\square$] Implement D3PM forward process
        (absorbing state, ternary space $\{0,1,m\}$).
  \item[$\square$] Implement 1D ResNet spectral encoder.
  \item[$\square$] Implement U-Net denoiser with AdaLN conditioning.
  \item[$\square$] Train with correct discrete CFG (Section~\ref{sec:cfg}).
  \item[$\square$] Validate unconditional generation quality
        (topology validity, diversity).
  \item[$\square$] Validate conditional generation
        (spectral accuracy under pure CFG, no physics guidance).
\end{itemize}

\subsection{Phase 4 — Physics-Guided Sampling (Weeks 12--17)}

\begin{itemize}
  \item[$\square$] Implement gradient guidance via predicted $\hat{x}_0$
        (Section~\ref{sec:guidance}).
  \item[$\square$] Implement uncertainty-weighted guidance.
  \item[$\square$] Implement connectivity discriminator and topology guidance.
  \item[$\square$] Implement differentiable DRC loss.
  \item[$\square$] Tune guidance hyperparameters
        ($\alphamax$, $w$, $\tthresh$, $\lambda$ weights).
  \item[$\square$] Validate: full-wave EM verification of generated layouts.
\end{itemize}

\subsection{Phase 5 — Evaluation and Paper (Weeks 15--22)}

\begin{itemize}
  \item[$\square$] Implement all baselines (Det-CNN, cVAE, cGAN, BO, Diff-TO).
  \item[$\square$] Run full evaluation on 1000 test specifications.
  \item[$\square$] Run all ablation studies.
  \item[$\square$] Perform scaling experiment at $32\times32$.
  \item[$\square$] Write paper (AAAI format, 8 pages + references).
  \item[$\square$] Prepare supplementary material
        (full architecture specs, training details, extended results).
\end{itemize}

% ============================================================
\section{Compute Budget and Resource Analysis}
\label{sec:compute}
% ============================================================

\subsection{Dataset Generation}

\begin{table}[H]
\centering
\caption{Dataset generation compute estimate (200k structures).}
\label{tab:compute_data}
\begin{tabular}{lll}
\toprule
\textbf{Parallelism} & \textbf{Wall Time} & \textbf{Cloud Cost (est.)} \\
\midrule
1 CPU core      & $\sim$6700~hours & --- \\
50 CPU cores    & $\sim$5.6~days   & --- \\
100 CPU cores   & $\sim$2.8~days   & --- \\
2$\times$ EC2 c5.18xlarge (72 vCPU each)
                & $\sim$2.8~days   & $\approx\$420$ \\
\bottomrule
\end{tabular}
\end{table}

\subsection{Surrogate Training}

\begin{itemize}
  \item Per surrogate: 100 epochs $\times$ $\sim$20~min/epoch on A100 $= \sim$33~hours.
  \item $K = 5$ surrogates in parallel on 5 GPUs: $\sim$33~hours $\approx$ 1.5~days.
  \item Estimated cost: \$$500$--\$$660$.
\end{itemize}

\subsection{Denoiser Training}

\begin{itemize}
  \item 300 epochs $\times$ $\sim$45~min/epoch on A100 $= \sim$225~hours $\approx$ 9~days.
  \item With 2 GPUs (data parallel): $\sim$4.5~days.
  \item Estimated cost: \$$700$--\$$900$.
\end{itemize}

\subsection{Total Budget Summary}

\begin{table}[H]
\centering
\caption{Total compute budget.}
\label{tab:compute_total}
\begin{tabular}{lll}
\toprule
\textbf{Phase} & \textbf{Est.\ Cost} & \textbf{Parallelised Wall Time} \\
\midrule
Dataset generation (200k)   & \$$400$--\$$600$   & $\sim$3--5~days \\
Surrogate training ($\times5$) & \$$500$--\$$700$ & $\sim$2--3~days \\
Denoiser training           & \$$700$--\$$900$   & $\sim$5--7~days \\
Evaluation + baselines      & \$$200$--\$$300$   & $\sim$2--3~days \\
\midrule
\textbf{Total}              & $\approx$\$$1800$--\$$2500$
                            & $\approx$12--18~days active compute \\
\bottomrule
\end{tabular}
\end{table}

% ============================================================
\section{Risk Analysis and Contingency Plans}
\label{sec:risks}
% ============================================================

\begin{table}[H]
\centering
\caption{Risk register with probability, impact, and mitigation.}
\label{tab:risks}
\begin{tabularx}{\linewidth}{lllX}
\toprule
\textbf{Risk} & \textbf{Prob.} & \textbf{Impact} & \textbf{Mitigation} \\
\midrule
R1: Poor surrogate gradient quality
  & Medium & High
  & Train on soft inputs; gradient smoothing;
    zero-order guidance fallback
    (finite-difference surrogate sampling). \\
\addlinespace
R2: D3PM ELBO training instability
  & Low & Medium
  & Switch to MDLM (masked diffusion) or simple
    Bernoulli diffusion as fallback. \\
\addlinespace
R3: Dataset quality insufficient
  & Low--Med & High
  & Include 10\% random-plus-repair structures;
    validate each sample with holdout EM simulation. \\
\addlinespace
R4: Compute budget overrun
  & Medium & Medium
  & Use AWS Spot instances;
    reduce target to 100k if 200k infeasible;
    use ML-accelerated FDTD for early dataset. \\
\addlinespace
R5: AAAI venue mismatch
  & Low--Med & Low
  & Ensure primary framing is AI methodology;
    prepare parallel submission to NeurIPS-2027
    or ICLR-2027. \\
\bottomrule
\end{tabularx}
\end{table}

% ============================================================
%  APPENDICES
% ============================================================
\appendix

% ============================================================
\section{Key Equations Reference}
\label{app:equations}
% ============================================================

\subsection{Physics}

\begin{table}[H]
\centering
\caption{Core physics equations.}
\label{tab:app_physics}
\begin{tabular}{ll}
\toprule
\textbf{Equation} & \textbf{Description} \\
\midrule
$\Smat^\dagger\Smat \preceq \mathbf{I}$ & Passivity ($N$-port) \\
$\Smat^T = \Smat$ & Reciprocity \\
$\Real[S] = \Hilbert\{\Imag[S]\}$ & Kramers--Kronig causality \\
$\Delta_{\mathrm{pixel}} \ll \lambda_{\mathrm{eff}}/10$ & Pixel size validity criterion \\
\bottomrule
\end{tabular}
\end{table}

\subsection{Dataset Parameters}

\begin{table}[H]
\centering
\caption{Dataset configuration.}
\label{tab:app_dataset}
\begin{tabular}{ll}
\toprule
\textbf{Parameter} & \textbf{Value} \\
\midrule
Grid size          & $15\times15$ (extendable to $32\times32$) \\
Frequency range    & $0.5$--$20$~GHz \\
Frequency points $N_f$ & 100 \\
Target size        & 100k--500k structures \\
Substrate types    & $\varepsilon_r \in \{2.2,\; 3.5,\; 4.4,\; 10.2\}$ \\
\bottomrule
\end{tabular}
\end{table}

\subsection{Model Hyperparameters}

\begin{table}[H]
\centering
\caption{Model hyperparameters.}
\label{tab:app_hyper}
\begin{tabular}{ll}
\toprule
\textbf{Parameter} & \textbf{Value} \\
\midrule
Spectral embedding dim $m$     & 128 or 256 \\
Diffusion steps $T$             & 500--1000 \\
CFG guidance weight $w$         & 2.0 (tunable) \\
Physics guidance $\alphamax$    & 0.10 \\
Guidance threshold $\tthresh$   & $0.4T$ \\
Ensemble size $K$               & 5 \\
Temperature initial $\tau_0$    & 1.0 \\
Temperature final $\tau_T$      & 0.01 \\
\bottomrule
\end{tabular}
\end{table}

% ============================================================
\section{Terminology Glossary}
\label{app:glossary}
% ============================================================

\begin{table}[H]
\centering
\caption{Glossary of terms used throughout this document.}
\label{tab:glossary}
\begin{tabularx}{\linewidth}{lX}
\toprule
\textbf{Term} & \textbf{Definition} \\
\midrule
S-parameters    & Scattering parameters describing network port behaviour \\
$S_{11}$        & Reflection coefficient at port~1 (return loss) \\
$S_{21}$        & Forward transmission from port~1 to port~2 (insertion loss) \\
FDTD            & Finite-Difference Time-Domain EM solver method \\
D3PM            & Discrete Denoising Diffusion Probabilistic Models (Austin et al., 2021) \\
CFG             & Classifier-Free Guidance \\
AdaLN           & Adaptive Layer Normalisation \\
STE             & Straight-Through Estimator \\
DRC             & Design Rule Check (manufacturing constraint verification) \\
KK              & Kramers--Kronig (causality relations between real and imaginary parts) \\
EM manifold     & The low-dimensional manifold of physically meaningful EM structures
                  within the full binary space \\
Procedural generation
                & Generating dataset samples from parameterised physical primitives
                  rather than random sampling \\
Absorbing state & A MASK token in D3PM that pixels transition into and cannot leave
                  in the forward process \\
ELBO            & Evidence Lower BOund — the training objective for variational models \\
GradNorm        & Gradient normalisation technique for multi-task learning
                  (Chen et al., 2018) \\
\bottomrule
\end{tabularx}
\end{table}

% ============================================================
\section{Document Version History}
\label{app:versions}
% ============================================================

\begin{table}[H]
\centering
\caption{Document change history.}
\label{tab:versions}
\begin{tabularx}{\linewidth}{llX}
\toprule
\textbf{Version} & \textbf{Date} & \textbf{Changes} \\
\midrule
v1.0 & May 2026 (early) & Initial brainstorm; full pipeline outlined. \\
v2.0 & May 2026         & Four critical mathematical errors corrected
                          (CFG, gradient guidance, KK, passivity);
                          resolution--frequency co-design added;
                          new baselines (cVAE, cGAN) added;
                          full vulnerability register added;
                          all equations verified for correctness. \\
\bottomrule
\end{tabularx}
\end{table}

\bigskip
\begin{center}
\rule{0.6\linewidth}{0.4pt}\\[4pt]
{\small\color{gray}
  \textit{Document Version: 2.0 \quad|\quad May 2026 \quad|\quad Status: Active Research Blueprint}\\[2pt]
  \textit{All V1 mathematical errors identified, documented, and corrected.}\\[2pt]
  \textit{Ready for implementation Phase~1 start.}\\[8pt]
  \textit{Hare Krishna. Hari Bol.}
}
\end{center}

\end{document}







